\chapter{状态空间建模与求解}
\label{ch:statespace-basics}

% ════════════════════════════════════════════════════════════════
%  P4 控制部「C0 现代控制理论基础」子部 · 第 C0-1 章
%  主源：刘豹《现代控制理论》(第3版) ch1（状态空间表达式）+ ch2（状态方程的解）
%  辅源/互校：张嗣瀛·高立群《现代控制理论》(第2版) ch2（m!=0 实现/组合系统/离散）+ ch3（时变 STM / Peano–Baker）
%  定位：把导论 \cref{ch:control} 的「给结论」状态空间记号，下沉为「会推 + 物理直觉 + 图解」的教科书地基；
%        向上为 A–Q 与 C2 动力学供给 (A,B) 与离散化形式。
%  教学层遵循 v5.0：动机→反面→历史→理论→陷阱→练习 + 前置自测/知识导航/小结速查。
%  记号归一：状态/输入粗体 x,u；机械系统 M(q)q̈+Cq̇+g=tau；引用教材公式前对照 OCR 勘误。
% ════════════════════════════════════════════════════════════════

\begin{practice}[前置自测]
开始之前，先试答下列问题。若已能流畅作答，可只速读本章的图与速查卡；若卡住，正是本章要为你解决的。
\begin{enumerate}
  \item 一个二阶 RLC 电路、一台直流电机、一根带弹簧阻尼的机械臂关节，凭什么都能写成同一形式 $\dot{\mathbf{x}}=\mathbf{A}\mathbf{x}+\mathbf{B}\mathbf{u}$？“状态”到底指什么、个数由什么决定？
  \item 同一个系统，你选不同的物理量当状态，得到的 $(\mathbf{A},\mathbf{B},\mathbf{C})$ 完全不同。那么哪些量是“跟选取无关”的系统固有属性？
  \item 标量方程 $\dot x=ax$ 的解是 $x(t)=\mathrm{e}^{at}x_0$。把 $a$ 换成矩阵 $\mathbf{A}$，“$\mathrm{e}^{\mathbf{A}t}$”是什么意思？它真的能像标量那样无脑套用吗（例如 $\mathrm{e}^{\mathbf{A}t}\mathrm{e}^{\mathbf{B}t}=\mathrm{e}^{(\mathbf{A}+\mathbf{B})t}$ 成立吗）？
  \item 给你 $\mathbf{A}$，你能说出至少三种算 $\mathrm{e}^{\mathbf{A}t}$ 的办法吗？它们各自的代价与适用场景是什么？
  \item 机器人沿一条轨迹做线性化，得到的是\emph{时变}系统 $\dot{\mathbf{x}}=\mathbf{A}(t)\mathbf{x}$。此时还能写 $\mathbf{x}(t)=\mathrm{e}^{\int_{t_0}^t\mathbf{A}(\tau)\mathrm{d}\tau}\mathbf{x}_0$ 吗？若不能，差在哪？
  \item 要把连续模型交给数字计算机或写成 MPC 的预测方程，得到 $\mathbf{x}_{k+1}=\mathbf{A}_d\mathbf{x}_k+\mathbf{B}_d\mathbf{u}_k$。$\mathbf{A}_d,\mathbf{B}_d$ 与连续的 $\mathbf{A},\mathbf{B}$ 是什么关系？
\end{enumerate}
\end{practice}

\section*{本章目标}
学完本章，你应当能够：
\begin{enumerate}
  \item \textbf{说清}“状态”的物理本质（最小的、足以确定未来的一组变量 $=$ 独立储能元件数），并辨析它与输入/输出的区别；
  \item \textbf{手把手列写}状态空间表达式的三条途径——结构图、物理机理、高阶微分方程/传递函数——并独立完成 RLC/机械/电机的统一示例链；
  \item \textbf{处理} $m\neq0$（输出含输入导数项）的实现，给出能控/能观标准形，并解释“为何不能直接拿输出各阶导数当状态”；
  \item \textbf{运用}非奇异线性变换 $\mathbf{x}=\mathbf{T}\mathbf{z}$，并\emph{证明}特征值与特征多项式系数是变换不变量（系统不变量）；
  \item \textbf{化}任意 $\mathbf{A}$ 为对角/约旦标准型，理解约旦块对应“耦合的模态”；
  \item \textbf{建立直觉}并\emph{推导}矩阵指数 $\mathrm{e}^{\mathbf{A}t}$（幂级数定义、几何轨线、组合/求逆性质），并掌握其四种算法（定义级数、约旦型、拉氏反变换、Cayley--Hamilton）；
  \item \textbf{求解}线性定常非齐次方程，把解读为“自由运动 $+$ 卷积强迫运动”；
  \item \textbf{推广}到\emph{时变}系统：理解可交换条件为何苛刻、用 Peano--Baker 级数与基本解阵定义状态转移矩阵 $\boldsymbol{\Phi}(t,t_0)$，为机器人沿轨迹线性化（LTV）打底；
  \item \textbf{离散化}连续模型（零阶保持精确式与一阶近似式），为数字控制、MPC 预测方程、ZOH-LIPM 供给 $(\mathbf{A}_d,\mathbf{B}_d)$。
\end{enumerate}

\subsection*{本章知识导航}
本章是一条“\emph{把物理系统写成 $(\mathbf{A},\mathbf{B},\mathbf{C},\mathbf{D})$，再把它在时间上解出来}”的主线。前半（建模）回答“系统长什么样”，后半（求解）回答“系统怎么动”。结构如下：

\begin{center}
\small
\begin{tikzpicture}[
  node distance=4.5mm and 7mm, >=Stealth,
  blk/.style={draw, rounded corners, align=center, font=\small, inner sep=3pt, fill=main!6, draw=main!60!black},
  grp/.style={draw, rounded corners, align=center, font=\small, inner sep=3pt, fill=second!6, draw=second!60!black},
  outb/.style={draw, rounded corners, align=center, font=\small, inner sep=3pt, fill=third!8, draw=third!60!black},
  ln/.style={->, thick, gray!60!black}]
\node[blk] (def) {状态 / 状态空间\\ $\dot{\mathbf{x}}=\mathbf{A}\mathbf{x}+\mathbf{B}\mathbf{u}$};
\node[blk, right=of def] (model) {三途径建模\\ 机理·结构图·传函};
\node[blk, right=of model] (trans) {线性变换\\ 不变量·约旦型};
\node[grp, below=8mm of def] (exp) {矩阵指数 $\mathrm{e}^{\mathbf{A}t}$\\ 直觉·性质};
\node[grp, right=of exp] (algo) {四算法\\ 级数/约旦/拉氏/C-H};
\node[grp, right=of algo] (nonhom) {非齐次解\\ 自由 $+$ 卷积};
\node[outb, below=8mm of exp] (ltv) {时变 STM\\ Peano-Baker $\boldsymbol{\Phi}(t,t_0)$};
\node[outb, right=of ltv] (disc) {离散化\\ $(\mathbf{A}_d,\mathbf{B}_d)$\\ $\to$ MPC/数字控制};
\draw[ln] (def)--(model); \draw[ln] (model)--(trans);
\draw[ln] (trans.south) -- ++(0,-4.5mm) -| (exp.north);
\draw[ln] (exp)--(algo); \draw[ln] (algo)--(nonhom);
\draw[ln] (exp.south)--(ltv.north);
\draw[ln] (nonhom.south) -- ++(0,-4.5mm) -| (disc.north);
\draw[ln] (ltv)--(disc);
\end{tikzpicture}
\end{center}

\noindent\textbf{推荐路径}：第 1--3 节是任何读者的主干（建立“状态—建模—不变量”三角），机器人/多体动力学读者尤应吃透机-电统一示例链与 $m\neq0$ 实现；第 4--6 节是\emph{求解}的核心工具（$\mathrm{e}^{\mathbf{A}t}$、四算法、非齐次解），后续能控能观、李雅普诺夫方程、LQR/MPC 都建立在 $\mathrm{e}^{\mathbf{A}t}$ 与离散化之上；第 7 节（时变）与第 8 节（离散化）面向工程落地——沿轨迹线性化得到 LTV、把模型交给数字处理器/MPC 时必经此关，初读可先掌握结论、用时回看。

\subsection*{与导论的分工}
导论 \cref{ch:control} 在 \cref{sec:ctrl-statespace} 已\emph{陈述}了状态空间四元组 $(\mathbf{A},\mathbf{B},\mathbf{C},\mathbf{D})$ 与机器人动力学的状态空间结构，并直接给出了可控/可观判据（\cref{def:ctrl-controllable}）、李雅普诺夫方程（\cref{eq:ctrl-lyap-eq}）、LQR/Riccati（\cref{sec:ctrl-lqr}）等\emph{结论}——那是“会用、给全景”的导论密度。本章与之分工明确：导论给\emph{会用}，本章给\emph{会推 $+$ 物理直觉}。凡导论“给结论”之处，本章用半句 \verb|\cref{}| 接住、随即下沉展开；\textbf{不重复编号同名定义/定理}，只把它们的来由、推导与几何图像补全。具体而言：本章把状态空间表达式从“一个矩阵方程”还原为“物理系统的最小动态描述”，把 $\mathrm{e}^{\mathbf{A}t}$ 从“记号”还原为“可推、可算、有几何意义的状态转移”，从而为整个 C0 子部（能控能观、稳定性、综合、最优控制基础）与下游的 A--Q 高级章夯实地基。

\subsection*{前置知识桥接}
\cref{ch:rigid_body} 与 \cref{ch:lie} 已建立刚体运动与流形上的记号（旋转 $\mathbf{R}$、位姿 $\mathbf{T}$、机械系统动力学 $\mathbf{M}(\mathbf{q})\ddot{\mathbf{q}}+\mathbf{C}\dot{\mathbf{q}}+\mathbf{g}=\boldsymbol{\tau}$）。本章建立的状态空间记号与之衔接：把二阶动力学“降阶”为一阶 $\dot{\mathbf{x}}=\mathbf{A}\mathbf{x}+\mathbf{B}\mathbf{u}$ 是后续 C2 动力学（$\mathbf{J}/\mathbf{K}/\mathbf{L}$）与浮动基座控制的统一入口。线性代数（特征值、相似变换、约旦型）与拉氏变换为假定背景，本章直接使用、不再重证。

\begin{note}
\textbf{本章记号与约定.}\;
状态矢量 $\mathbf{x}\in\mathbb{R}^n$、输入（控制）$\mathbf{u}\in\mathbb{R}^r$、输出 $\mathbf{y}\in\mathbb{R}^m$，一律\textbf{粗体}。系统四元组：系统矩阵 $\mathbf{A}\in\mathbb{R}^{n\times n}$、输入矩阵 $\mathbf{B}\in\mathbb{R}^{n\times r}$、输出矩阵 $\mathbf{C}\in\mathbb{R}^{m\times n}$、直接传递矩阵 $\mathbf{D}\in\mathbb{R}^{m\times r}$。单输入单输出（SISO）时用小写 $\mathbf{b},\mathbf{c},d$。状态转移矩阵记 $\boldsymbol{\Phi}(t)=\mathrm{e}^{\mathbf{A}t}$（定常）、$\boldsymbol{\Phi}(t,t_0)$（时变）。
\textbf{记号归一}：两本源教材用非粗体、约旦型记 $\mathbf{J}$、传函记 $W(s)$、增益记 $K_P/K_D$——本书融入时统一为上述粗体记号；引用教材经典结论处 \verb|\cite{}| 标源。
\textbf{P/Q/R 双关纪律}（沿用 C 章规矩、与 \cref{ch:eskf} 区分）：控制侧 $\mathbf{P}$=代价权重/Riccati 解、$\mathbf{Q},\mathbf{R}$=代价权重；估计侧噪声协方差 $\boldsymbol{\Sigma}_w,\boldsymbol{\Sigma}_v$。本章只涉及前者。
\end{note}

% ════════════════════════════════════════════════════════════════
\section{从物理系统到状态空间表达式}
\label{sec:ssb-state}

\paragraph{动机：经典描述“看不全”系统。}
经典控制理论用一个输出量的高阶微分方程或传递函数描述系统——它只把\emph{某个}输出直接和输入挂钩\cite{liubao2006modern}。但一个系统内部还有许多\emph{相互独立}的中间变量，微分方程/传递函数对它们不便描述，因而不能包含系统的全部信息。从“能否完全揭示系统全部运动”来看，这种外部描述有其根本不足。状态空间法换一个思路：用一组\emph{一阶}微分方程同时刻画系统的全部独立变量，从而完整确定系统的内部运动，还能方便地处理初始条件、统一对待时变/非线性/多输入多输出/随机系统——这正是现代控制（最优控制、状态反馈、卡尔曼滤波）的共同地基。

\paragraph{核心概念：什么是“状态”。}
\begin{definition}[状态变量与状态矢量]\label{def:ssb-state}
\textbf{状态变量}是足以完全表征系统运动、而个数又\emph{最小}的一组变量 $x_1(t),\dots,x_n(t)$：当它们在初始时刻 $t_0$ 的值 $\mathbf{x}(t_0)$ 已知、且 $t\ge t_0$ 的输入给定时，便能唯一确定系统在任意 $t\ge t_0$ 的行为。把它们叠成列向量
\begin{equation}
  \mathbf{x}(t)=\big[x_1(t),\,x_2(t),\,\dots,\,x_n(t)\big]^\top\in\mathbb{R}^n
  \label{eq:ssb-statevec}
\end{equation}
即\textbf{状态矢量}。以 $x_1,\dots,x_n$ 为坐标轴张成的 $n$ 维空间称\textbf{状态空间}；$\mathbf{x}(t)$ 在其中是一点，随时间画出一条\textbf{状态轨线}。
\end{definition}

\begin{insight}
“状态个数 $=$ 独立储能元件数”是最有用的物理直觉\cite{liubao2006modern}。$n$ 阶微分方程要唯一定解须知 $n$ 个独立初始条件，而微分方程的阶数恰由系统中\emph{独立储能元件}个数决定（电容储电压能、电感储电流能、质量储动能、弹簧储势能）。所以：数一数系统里有几个独立的“能量蓄水池”，就知道要几个状态变量。RLC 电路里电容电压 $u_C$ 与电感电流 $i$、机械系统里位移与速度，都是天然的状态——它们是“记住过去、决定未来”的最小信息。
\end{insight}

\paragraph{招牌示例：RLC 电路的状态方程。}
取图示串联 RLC 电路（电感 $L$、电阻 $R$、电容 $C$、电源 $u$），有两个独立储能元件，故两个状态。以电容电压 $u_C$、电感电流 $i$ 为状态（它们直接对应储能）。由基尔霍夫定律：
\begin{align}
  C\,\dot{u}_C &= i, &
  L\,\dot{i} &= -u_C - R\,i + u.
  \label{eq:ssb-rlc-phys}
\end{align}
\begin{note}
\textbf{OCR 勘误（刘豹式1.1处）.}\;源 OCR 把电容电压下标 $C$ 后误并入一个句点（$u_{c.}$），输出方程处 $u_C$ 被识成 $u_{cc}$；本书已据勘误更正为 $u_C$\cite{liubao2006modern}。
\end{note}
令 $x_1=u_C,\,x_2=i$，写成矩阵形式即得状态方程，再指定 $y=u_C$ 为输出得输出方程：
\begin{equation}
  \underbrace{\begin{bmatrix}\dot{x}_1\\\dot{x}_2\end{bmatrix}}_{\dot{\mathbf{x}}}
  =\underbrace{\begin{bmatrix}0&\tfrac{1}{C}\\[2pt]-\tfrac{1}{L}&-\tfrac{R}{L}\end{bmatrix}}_{\mathbf{A}}
  \begin{bmatrix}x_1\\x_2\end{bmatrix}
  +\underbrace{\begin{bmatrix}0\\\tfrac{1}{L}\end{bmatrix}}_{\mathbf{b}}u,
  \qquad
  y=\underbrace{\begin{bmatrix}1&0\end{bmatrix}}_{\mathbf{c}}\mathbf{x}.
  \label{eq:ssb-rlc-ss}
\end{equation}

\paragraph{形式化：一般状态空间表达式。}
\begin{definition}[线性定常系统的状态空间表达式]\label{def:ssb-ss}
\emph{状态方程}（由状态变量构成的一阶微分方程组）与\emph{输出方程}（输出与状态/输入的代数关系）合起来，构成系统完整的动态描述。多输入（$r$ 维）多输出（$m$ 维）线性定常系统为
\begin{equation}
  \boxed{\;\dot{\mathbf{x}}=\mathbf{A}\mathbf{x}+\mathbf{B}\mathbf{u},\qquad \mathbf{y}=\mathbf{C}\mathbf{x}+\mathbf{D}\mathbf{u}\;}
  \label{eq:ssb-mimo}
\end{equation}
其中 $\mathbf{A}_{n\times n}$ 刻画\emph{内部状态间的耦合}（系统矩阵），$\mathbf{B}_{n\times r}$ 刻画\emph{输入对状态的作用}（输入/控制矩阵），$\mathbf{C}_{m\times n}$ 刻画\emph{状态对输出的影响}（输出矩阵），$\mathbf{D}_{m\times r}$ 是输入到输出的\emph{直接传递}（一般为零）。
\end{definition}

这与导论 \cref{sec:ctrl-statespace} 写出的四元组完全一致；导论从“MIMO 控制设计的入口”引入它，本章则强调它是物理系统的\emph{最小完整描述}。其框图把因果关系画得很清楚：积分器数 $=$ 状态数，每个积分器输出一个 $x_i$、输入是 $\dot x_i$，再用加法器/比例器按 $\mathbf{A},\mathbf{B},\mathbf{C}$ 连线——单线表标量、双线表矢量。

\begin{figure}[ht]
\centering
\begin{tikzpicture}[>=Stealth, node distance=10mm, every node/.style={font=\small},
  blk/.style={draw, rounded corners, minimum height=7mm, minimum width=10mm, fill=main!6, draw=main!60!black},
  sum/.style={draw, circle, inner sep=1.2pt, fill=second!10, draw=second!60!black}]
% MIMO vector block diagram (double-line = vector). Drawn purely with explicit coords.
\node (u) at (0,0) {$\mathbf{u}$};
\node[sum] (s1) at (1.6,0) {$+$};
\node[blk] (B) at (1.6,1.1) {$\mathbf{B}$};
\node[blk] (int) at (3.6,0) {$\displaystyle\int$};
\node (xdot) at (2.65,0.32) {$\dot{\mathbf{x}}$};
\node (x) at (4.7,0.32) {$\mathbf{x}$};
\node[blk] (C) at (6.1,0) {$\mathbf{C}$};
\node[sum] (s2) at (7.5,0) {$+$};
\node[blk] (D) at (4.7,1.6) {$\mathbf{D}$};
\node[blk] (A) at (3.6,-1.5) {$\mathbf{A}$};
\node (y) at (8.6,0) {$\mathbf{y}$};
\draw[->, very thick, double] (u)--(s1);
\draw[->, very thick, double] (u.east)++(0.3,0) |- (B.west);
\draw[->, very thick, double] (B)--(s1);
\draw[->, very thick, double] (s1)--(int);
\draw[->, very thick, double] (int)-- node[above]{} (C);
\draw[->, very thick, double] (4.2,0)|-(A.east);
\draw[->, very thick, double] (A.west)-|(s1.south);
\draw[->, very thick, double] (C)--(s2);
\draw[->, very thick, double] (u.east)++(0.3,0)|-(D.west);
\draw[->, very thick, double] (D)-|(s2.north);
\draw[->, very thick, double] (s2)--(y);
\end{tikzpicture}
\caption{多输入多输出系统 \cref{eq:ssb-mimo} 的矢量框图：积分器吐出状态 $\mathbf{x}$，其输入 $\dot{\mathbf{x}}$ 由 $\mathbf{A}\mathbf{x}+\mathbf{B}\mathbf{u}$ 汇成；双线箭头表矢量信号。状态反馈 $\mathbf{A}$ 是回路，$\mathbf{D}$ 是前馈直通。（仿 \cite{liubao2006modern} 重绘）}
\label{fig:ssb-blockdiagram}
\end{figure}

\begin{pitfall}[概念误区：状态空间表达式唯一]
同一个系统，状态变量\emph{选取不唯一}，因而 $(\mathbf{A},\mathbf{B},\mathbf{C},\mathbf{D})$ 也\emph{不唯一}。例如 \cref{eq:ssb-rlc-ss} 选 $(u_C,i)$；若改选 $(u_C,\dot u_C)$ 会得到形式不同的另一组矩阵。这不是错误，而是状态空间法的\emph{自由度}——后面 \cref{sec:ssb-transform} 会证明：不论怎么选，特征值这样的“系统固有属性”始终不变。把“矩阵变了”误当“系统变了”，是初学者最常见的混淆。
\end{pitfall}

\begin{practice}
\begin{enumerate}
  \item 对 \cref{eq:ssb-rlc-ss}，从状态方程消去 $i$，验证以 $u_C$ 为输出的二阶微分方程为 $\ddot u_C+\tfrac{R}{L}\dot u_C+\tfrac{1}{LC}u_C=\tfrac{1}{LC}u$，并写出其传递函数。
  \item 对同一 RLC 电路改选 $x_1=u_C,\,x_2=\dot u_C$，重新列写 $(\mathbf{A},\mathbf{b})$，并与 \cref{eq:ssb-rlc-ss} 对比，体会“非唯一性”。
  \item 数一数：一个含 2 个电容、2 个电感的无源网络，状态空间维数是多少？
\end{enumerate}
\end{practice}

% ════════════════════════════════════════════════════════════════
\section{建立状态空间表达式的三条途径}
\label{sec:ssb-build}

建立状态空间表达式一般有三条途径\cite{liubao2006modern}：\emph{(i) 从系统结构图}（按各环节模拟结构图直接写）、\emph{(ii) 从物理机理}（用基尔霍夫/牛顿/能量守恒等物理定律列写）、\emph{(iii) 从高阶微分方程或传递函数}（“实现”问题）。本节用一条贯穿电、机、机电的\textbf{统一示例链}打磨途径 (ii)，再系统讲途径 (iii)——尤其是输出含输入导数（$m\neq0$）的处理，这是机器人传感器动力学常遇的情形。

\subsection{机理途径：机-电统一示例链}
\label{sec:ssb-mechanism}

按物理规律列写最直观：选独立储能元件的物理量为状态，写出各储能元件的动态方程，整理成 \cref{eq:ssb-mimo} 即可。下面三例同出一辙，凸显“不同物理、同一形式”。

\paragraph{例 1（机械平动）：弹簧-质量-阻尼。}
两质量块 $M_1,M_2$、弹簧 $K_1,K_2$、阻尼 $B_1,B_2$，外力 $f$ 作用、以位移 $y_1,y_2$ 为输出。弹簧与质量是储能元件，故选位移与速度为状态：$x_1=y_1,\,x_2=y_2,\,x_3=\dot y_1,\,x_4=\dot y_2$。对 $M_1,M_2$ 用牛顿第二定律
\begin{align}
  M_1\ddot y_1 &= K_2(y_2-y_1)+B_2(\dot y_2-\dot y_1)-K_1 y_1-B_1\dot y_1, \\
  M_2\ddot y_2 &= f-K_2(y_2-y_1)-B_2(\dot y_2-\dot y_1),
  \label{eq:ssb-mech-newton}
\end{align}
代入状态变量整理（$u=f$）得
\begin{equation}
  \dot{\mathbf{x}}=
  \begin{bmatrix}
    0&0&1&0\\ 0&0&0&1\\
    -\tfrac{K_1+K_2}{M_1}&\tfrac{K_2}{M_1}&-\tfrac{B_1+B_2}{M_1}&\tfrac{B_2}{M_1}\\[2pt]
    \tfrac{K_2}{M_2}&-\tfrac{K_2}{M_2}&\tfrac{B_2}{M_2}&-\tfrac{B_2}{M_2}
  \end{bmatrix}\mathbf{x}
  +\begin{bmatrix}0\\0\\0\\\tfrac{1}{M_2}\end{bmatrix}u,
  \quad
  \mathbf{y}=\begin{bmatrix}1&0&0&0\\0&1&0&0\end{bmatrix}\mathbf{x}.
  \label{eq:ssb-mech-ss}
\end{equation}

\begin{insight}
\cref{eq:ssb-mech-ss} 的上半 $\begin{bmatrix}\mathbf{0}&\mathbf{I}\\ \cdot&\cdot\end{bmatrix}$ 结构不是巧合：它正是机械系统“位置的导数是速度”的体现。把 \cref{ch:rigid_body} 的二阶动力学 $\mathbf{M}(\mathbf{q})\ddot{\mathbf{q}}+\mathbf{C}\dot{\mathbf{q}}+\mathbf{g}=\boldsymbol{\tau}$ 取状态 $\mathbf{x}=[\mathbf{q}^\top,\dot{\mathbf{q}}^\top]^\top$，就得到统一的“降阶”形式
\begin{equation}
  \dot{\mathbf{x}}=\begin{bmatrix}\dot{\mathbf{q}}\\\ddot{\mathbf{q}}\end{bmatrix}
  =\begin{bmatrix}\mathbf{0}&\mathbf{I}\\ -\mathbf{M}^{-1}\mathbf{K}&-\mathbf{M}^{-1}\mathbf{C}\end{bmatrix}\mathbf{x}
  +\begin{bmatrix}\mathbf{0}\\\mathbf{M}^{-1}\end{bmatrix}\boldsymbol{\tau}
  +\begin{bmatrix}\mathbf{0}\\-\mathbf{M}^{-1}\mathbf{g}\end{bmatrix},
  \label{eq:ssb-mechanical-general}
\end{equation}
（线性化时 $\mathbf{g}$ 项并入常量/仿射项）。这条“二阶动力学 $\to$ 一阶状态空间”是机械臂、足式、浮动基座控制的统一入口，下游 C2 动力学章（$\mathbf{J}/\mathbf{K}/\mathbf{L}$）正是从这里取用 $(\mathbf{A},\mathbf{B})$。
\end{insight}

\paragraph{例 2（机械转动）：扭转轴。}
转动惯量 $J$、扭转刚度 $K$、粘滞阻尼 $B$、力矩 $T$。选 $x_1=\theta,\,x_2=\dot\theta,\,u=T$，由 $J\ddot\theta=-K\theta-B\dot\theta+T$ 得
\begin{equation}
  \dot{\mathbf{x}}=\begin{bmatrix}0&1\\-\tfrac{K}{J}&-\tfrac{B}{J}\end{bmatrix}\mathbf{x}+\begin{bmatrix}0\\\tfrac{1}{J}\end{bmatrix}u,
  \qquad y=\begin{bmatrix}1&0\end{bmatrix}\mathbf{x}.
  \label{eq:ssb-rot-ss}
\end{equation}
这正是单关节的最简模型，与 \cref{eq:ssb-rlc-ss} 的二维结构如出一辙。

\paragraph{例 3（机电耦合）：直流他励电动机。}
电枢电阻 $R$、电感 $L$、转动惯量 $J$、粘滞摩擦 $B$；电流 $i$ 与转速 $\omega$ 为独立状态。电枢回路 $L\dot i+Ri+e=u$、动力学 $J\dot\omega+B\omega=K_a i$、反电动势 $e=K_b\omega$（$K_a$ 转矩常数、$K_b$ 反电动势常数）。消去 $e$、取 $x_1=i,x_2=\omega$：
\begin{equation}
  \dot{\mathbf{x}}=\begin{bmatrix}-\tfrac{R}{L}&-\tfrac{K_b}{L}\\[2pt]\tfrac{K_a}{J}&-\tfrac{B}{J}\end{bmatrix}\mathbf{x}+\begin{bmatrix}\tfrac{1}{L}\\0\end{bmatrix}u.
  \label{eq:ssb-motor-ss}
\end{equation}
若改以电机\emph{转角} $\theta$ 为输出，单凭 $(i,\omega)$ 不足以描述时域行为，须\emph{增}一个状态 $x_3=\theta$（$\dot x_3=\omega$），状态维数升到 3——这生动说明：状态个数由“要完整描述什么”决定。

\begin{insight}
三例放在一起，状态空间法的威力一目了然：电路（KVL/KCL）、平动（牛顿）、转动（转动定律）、机电（电路 $+$ 动力学 $+$ 电磁耦合）——物理机理千差万别，却都被收编为同一个 $\dot{\mathbf{x}}=\mathbf{A}\mathbf{x}+\mathbf{B}\mathbf{u}$。\emph{一套求解与设计工具（$\mathrm{e}^{\mathbf{A}t}$、能控能观、LQR）因而能一视同仁地用于所有这些系统}。这正是“现代控制理论”相对“经典（按对象分门别类）”的根本进步。
\end{insight}

\subsection{实现途径：从高阶微分方程到状态空间}
\label{sec:ssb-realization}

逆问题——由输入输出的高阶微分方程或传递函数反求状态空间表达式——称\textbf{实现（realization）}问题。考虑单变量 $n$ 阶系统
\begin{equation}
  y^{(n)}+a_{n-1}y^{(n-1)}+\dots+a_1\dot y+a_0 y=b_m u^{(m)}+\dots+b_1\dot u+b_0 u,\quad m\le n,
  \label{eq:ssb-highorder}
\end{equation}
对应传递函数 $W(s)=\dfrac{b_m s^m+\dots+b_0}{s^n+a_{n-1}s^{n-1}+\dots+a_0}$。

\begin{note}
\textbf{实现的存在与（非）唯一.}\;实现存在 $\iff m\le n$。$m<n$ 时 $\mathbf{D}=0$；$m=n$ 时 $\mathbf{D}=b_m\neq0$（输出含与输入直通的项）。实现\emph{非唯一}（状态选法无穷多），但只要分子分母无公因子（无零极点对消），$n$ 阶系统必有 $n$ 个独立状态、各实现的特征根相同——这种无对消的实现称\textbf{最小实现}\cite{liubao2006modern,zhang2017modern}。最小实现的严格刻画（能控且能观）留待 \cref{ch:controllability-obsv}。
\end{note}

\paragraph{无零点（$m=0$）的实现：相变量与友矩阵。}
当 $W(s)=\dfrac{b_0}{s^n+a_{n-1}s^{n-1}+\dots+a_0}$，取每个积分器输出为状态（即输出的各阶导数，称\emph{相变量}）$x_1=y/b_0,\ x_2=\dot x_1,\dots$，直接得
\begin{equation}
  \dot{\mathbf{x}}=\underbrace{\begin{bmatrix}
    0&1&0&\cdots&0\\ 0&0&1&\cdots&0\\
    \vdots&\vdots&\vdots&\ddots&\vdots\\ 0&0&0&\cdots&1\\
    -a_0&-a_1&-a_2&\cdots&-a_{n-1}
  \end{bmatrix}}_{\text{友矩阵 }\mathbf{A}}\mathbf{x}
  +\begin{bmatrix}0\\0\\\vdots\\0\\1\end{bmatrix}u,
  \qquad y=\begin{bmatrix}b_0&0&\cdots&0\end{bmatrix}\mathbf{x}.
  \label{eq:ssb-companion}
\end{equation}
此 $\mathbf{A}$ 称\textbf{友矩阵（companion matrix）}：主对角线上方全为 $1$，最后一行是特征多项式系数取负，其余为零。它将贯穿全书（能控标准形、极点配置都用它）。

\begin{example}[$m=0$ 实现]\label{ex:ssb-companion}
$\dddot y+6\ddot y+41\dot y+7y=6u$（三阶系统）。
\begin{note}\textbf{OCR 勘误.}\;源 OCR 把首项三阶导 $\dddot y$ 系统性识成 $\ddot y$（导数点丢失），状态方程左端的导数点亦丢失；本书据勘误恢复\cite{liubao2006modern}。\end{note}
取 $x_1=y/6,\,x_2=\dot y/6,\,x_3=\ddot y/6$，得 $\dot x_1=x_2,\ \dot x_2=x_3,\ \dot x_3=-7x_1-41x_2-6x_3+u$，即
\[
  \dot{\mathbf{x}}=\begin{bmatrix}0&1&0\\0&0&1\\-7&-41&-6\end{bmatrix}\mathbf{x}+\begin{bmatrix}0\\0\\1\end{bmatrix}u,\qquad y=\begin{bmatrix}6&0&0\end{bmatrix}\mathbf{x}.
\]
\end{example}

\paragraph{含输入导数（$m\neq0$）：为何不能硬选 $y,\dot y,\dots$ 为状态。}
若 $m\neq0$（右端含 $\dot u,\ddot u,\dots$），仍硬取 $y,\dot y,\dots,y^{(n-1)}$ 为状态，则最后一个状态方程会\emph{显含 $u$ 的导数项}\cite{zhang2017modern}：
\begin{equation}
  \dot x_n=-a_n x_1-\dots-a_1 x_n+\big(b_0 u^{(n)}+b_1 u^{(n-1)}+\dots+b_n u\big).
  \label{eq:ssb-bad-state}
\end{equation}
这会让状态在输入跳变时出现\emph{无穷大跳变}，解的存在唯一性被破坏——这是我们不想要的。补救之道是\emph{重新选状态}，把输入导数“吸收”进状态定义，使状态方程不含 $\dot u$。下面给出两种标准形（此处用 $a_1,\dots,a_n$ 记分母系数、$b_0,\dots,b_m$ 记分子系数，与 \cref{ch:control}/\cref{ch:controllability-obsv} 的能控能观标准形衔接）。

\begin{derivation}[$m<n$：能控标准形与能观标准形]
\label{der:ssb-canonical}
引入算子 $p=\mathrm{d}/\mathrm{d}t$，把 \cref{eq:ssb-highorder}（$m<n$，记 $y=\dfrac{b_0 p^m+\dots+b_m}{p^n+a_1 p^{n-1}+\dots+a_n}u$）分解为中间变量 $\tilde y$：
\[
  \tilde y^{(n)}+a_1\tilde y^{(n-1)}+\dots+a_n\tilde y=u,\qquad y=b_0\tilde y^{(m)}+\dots+b_m\tilde y.
\]
\textbf{能控标准形}：选 $x_1=\tilde y,\,x_2=\dot{\tilde y},\dots,x_n=\tilde y^{(n-1)}$，则状态方程是友矩阵形、输出方程把系数搬到 $\mathbf{C}$：
\begin{equation}
  \dot{\mathbf{x}}=\begin{bmatrix}0&1&\cdots&0\\ \vdots&\vdots&\ddots&\vdots\\ 0&0&\cdots&1\\ -a_n&-a_{n-1}&\cdots&-a_1\end{bmatrix}\mathbf{x}+\begin{bmatrix}0\\\vdots\\0\\1\end{bmatrix}u,\quad
  y=\begin{bmatrix}b_m&\cdots&b_0&0&\cdots&0\end{bmatrix}\mathbf{x}.
  \label{eq:ssb-controllable-canon}
\end{equation}
\textbf{能观标准形}：选 $x_n=y,\,x_{n-1}=\dot y+a_1 y,\dots$（把输出与其导数、输入逐层吸收，使状态方程不含 $\dot u$），得到 \cref{eq:ssb-controllable-canon} 的“转置对偶”：
\begin{equation}
  \dot{\mathbf{x}}=\begin{bmatrix}0&0&\cdots&-a_n\\ 1&0&\cdots&-a_{n-1}\\ 0&1&\cdots&-a_{n-2}\\ \vdots&\vdots&\ddots&\vdots\\ 0&0&\cdots&-a_1\end{bmatrix}\mathbf{x}+\begin{bmatrix}b_m\\ b_{m-1}\\ \vdots\\ b_0\\ \vdots\\0\end{bmatrix}u,\quad y=\begin{bmatrix}0&\cdots&0&1\end{bmatrix}\mathbf{x}.
  \label{eq:ssb-observable-canon}
\end{equation}
两者的系统矩阵互为转置、$\mathbf{B}_c^\top=\mathbf{C}_o$、$\mathbf{C}_c^\top=\mathbf{B}_o$——这就是\emph{能控-能观对偶}的雏形\cite{zhang2017modern}，\cref{ch:controllability-obsv} 会把它讲透。
\end{derivation}

\begin{note}
\textbf{$m=n$ 的情形.}\;先把真有理分式“严格真化”：$W(s)=b_0+\dfrac{(b_1-b_0a_1)p^{n-1}+\dots+(b_n-b_0 a_n)}{p^n+a_1 p^{n-1}+\dots+a_n}$，常数项 $b_0$ 即直接传递 $\mathbf{D}=b_0$，余下严格真部分按 \cref{der:ssb-canonical} 实现。例：$y^{(3)}+16y^{(2)}+194y^{(1)}+640y=4u^{(3)}+160u^{(1)}+720u$，$n=m=3$，算出 $b_3-b_0a_3=720-4\cdot640=-1840$ 等，得 $\mathbf{C}=[-1840,-616,-64]$、$\mathbf{D}=4$\cite{zhang2017modern}。
\end{note}

\begin{insight}
“为何不能直接拿输出的各阶导数当状态”是 $m\neq0$ 的关键认知\cite{zhang2017modern}。物理上，含 $\dot u$ 意味着输入的\emph{突变率}直接进入状态变化率——一个理想阶跃输入其导数是冲激，会把状态瞬间“踢”到无穷。机器人里这并非纯学术：带微分环节的传感器/作动器、含零点的伺服回路都属此列。正确做法是把“输入的瞬时贡献”预先吸收进状态定义，让状态保持有界、解保持良态。这也预告了一个工程戒律：\emph{慎用纯微分器}（噪声会被放大），后续观测器（\cref{ch:linear-synthesis}）正是为绕开微分而生。
\end{insight}

% ════════════════════════════════════════════════════════════════
\section{线性变换、系统不变量与约旦标准型}
\label{sec:ssb-transform}

上一节反复出现“状态选取不唯一”。本节把这种自由度形式化为\emph{非奇异线性变换}，回答关键问题：哪些量随变换而变、哪些是系统的\emph{固有不变量}？并给出最有用的一种变换——化为对角/约旦标准型，它把耦合系统解耦成（近乎）独立的模态，是 $\mathrm{e}^{\mathbf{A}t}$ 计算、能控能观结构分析的共同利器。

\subsection{线性变换与系统不变量}
\label{sec:ssb-invariant}

\paragraph{坐标变换。}
给定系统 \cref{eq:ssb-mimo}，任取非奇异矩阵 $\mathbf{T}$，令 $\mathbf{x}=\mathbf{T}\mathbf{z}$（即 $\mathbf{z}=\mathbf{T}^{-1}\mathbf{x}$，新状态是旧状态的线性组合），代入得新表达式
\begin{equation}
  \dot{\mathbf{z}}=\mathbf{T}^{-1}\mathbf{A}\mathbf{T}\,\mathbf{z}+\mathbf{T}^{-1}\mathbf{B}\,\mathbf{u},\qquad \mathbf{y}=\mathbf{C}\mathbf{T}\,\mathbf{z}+\mathbf{D}\mathbf{u}.
  \label{eq:ssb-similarity}
\end{equation}
$\mathbf{T}$ 任意非奇异，故 $(\mathbf{A},\mathbf{B},\mathbf{C})\mapsto(\mathbf{T}^{-1}\mathbf{A}\mathbf{T},\mathbf{T}^{-1}\mathbf{B},\mathbf{C}\mathbf{T})$ 给出无穷多“同一系统的不同坐标表示”。

\begin{theorem}[特征值与系统不变量在线性变换下不变]\label{thm:ssb-invariant}
非奇异变换 \cref{eq:ssb-similarity} 不改变系统矩阵的特征值；从而特征多项式 $|\lambda\mathbf{I}-\mathbf{A}|=\lambda^n+a_{n-1}\lambda^{n-1}+\dots+a_0$ 的全部系数 $a_{n-1},\dots,a_0$ 也不变。这些系数称\textbf{系统不变量}。
\end{theorem}
\begin{proof}
直接计算变换后的特征多项式，利用 $\lambda\mathbf{I}=\mathbf{T}^{-1}\lambda\mathbf{T}$ 与行列式乘性：
\[
  |\lambda\mathbf{I}-\mathbf{T}^{-1}\mathbf{A}\mathbf{T}|
  =|\mathbf{T}^{-1}(\lambda\mathbf{I}-\mathbf{A})\mathbf{T}|
  =|\mathbf{T}^{-1}|\,|\lambda\mathbf{I}-\mathbf{A}|\,|\mathbf{T}|
  =|\mathbf{T}^{-1}\mathbf{T}|\,|\lambda\mathbf{I}-\mathbf{A}|
  =|\lambda\mathbf{I}-\mathbf{A}|.
\]
故两者特征多项式恒等，特征值相同；特征值由系数唯一确定，反之系数（作为根的对称函数）也由特征值唯一确定，故系数不变。
\end{proof}

\begin{insight}
\cref{thm:ssb-invariant} 是“矩阵变了系统没变”的精确表述：特征值（系统的极点、模态频率与衰减率）是\emph{物理实在}，不依赖你用什么坐标记账。传递函数同样是变换不变量（\cref{sec:ssb-tf} 将见 $\mathbf{C}(s\mathbf{I}-\mathbf{A})^{-1}\mathbf{B}+\mathbf{D}$ 不随 $\mathbf{T}$ 变）。所以做控制分析时，\emph{先挑一个让结构最清楚的坐标}（如约旦型）是完全合法的——这正是下一小节的动机。
\end{insight}

\paragraph{特征矢量。}
若非零向量 $\mathbf{p}_i$ 经 $\mathbf{A}$ 作用只改变长度、不改变方向，即 $\mathbf{A}\mathbf{p}_i=\lambda_i\mathbf{p}_i$，则 $\mathbf{p}_i$ 是 $\mathbf{A}$ 对应特征值 $\lambda_i$ 的\textbf{特征矢量}。它是构造对角化变换矩阵的砖块。

\subsection{对角化与约旦标准型}
\label{sec:ssb-jordan}

\begin{theorem}[特征值互异时可对角化]\label{thm:ssb-diag}
设 $\mathbf{A}$ 有 $n$ 个互异特征值 $\lambda_1,\dots,\lambda_n$，对应特征矢量 $\mathbf{p}_1,\dots,\mathbf{p}_n$（线性无关）。以它们为列构成 $\mathbf{T}=[\mathbf{p}_1,\dots,\mathbf{p}_n]$，则
\begin{equation}
  \mathbf{T}^{-1}\mathbf{A}\mathbf{T}=\boldsymbol{\Lambda}=\mathrm{diag}(\lambda_1,\dots,\lambda_n).
  \label{eq:ssb-diagonalize}
\end{equation}
\end{theorem}
\begin{proof}
特征值互异保证 $\mathbf{p}_1,\dots,\mathbf{p}_n$ 线性无关，故 $\mathbf{T}$ 非奇异。由 $\mathbf{A}\mathbf{p}_i=\lambda_i\mathbf{p}_i$：
\[
  \mathbf{A}\mathbf{T}=\mathbf{A}[\mathbf{p}_1,\dots,\mathbf{p}_n]=[\lambda_1\mathbf{p}_1,\dots,\lambda_n\mathbf{p}_n]=[\mathbf{p}_1,\dots,\mathbf{p}_n]\,\mathrm{diag}(\lambda_1,\dots,\lambda_n)=\mathbf{T}\boldsymbol{\Lambda}.
\]
左乘 $\mathbf{T}^{-1}$ 即得 \cref{eq:ssb-diagonalize}。
\begin{note}\textbf{OCR 勘误.}\;源 OCR 在该证明的 $\mathbf{A}\mathbf{T}=(\lambda_1\mathbf{p}_1,\lambda_i\mathbf{p}_2,\dots)$ 处把第 $2,\dots,n$ 列的下标误识成 $i$；正确为 $\lambda_k\mathbf{p}_k$\cite{liubao2006modern}。\end{note}
\end{proof}

\begin{example}[对角化]\label{ex:ssb-diag}
$\mathbf{A}=\begin{bmatrix}0&1&-1\\-6&-11&6\\-6&-11&5\end{bmatrix}$，特征方程 $\lambda^3+6\lambda^2+11\lambda+6=(\lambda+1)(\lambda+2)(\lambda+3)=0$，得 $\lambda_1=-1,\lambda_2=-2,\lambda_3=-3$。解 $\mathbf{A}\mathbf{p}_i=\lambda_i\mathbf{p}_i$ 得特征矢量
$\mathbf{p}_1=[1,0,1]^\top,\ \mathbf{p}_2=[1,2,4]^\top,\ \mathbf{p}_3=[1,6,9]^\top$，于是 $\mathbf{T}=[\mathbf{p}_1\,\mathbf{p}_2\,\mathbf{p}_3]$ 把 $\mathbf{A}$ 对角化为 $\mathrm{diag}(-1,-2,-3)$。
\end{example}

\paragraph{重根：约旦块与广义特征矢量。}
特征值有重根且独立特征矢量不足时，不能完全对角化，只能化为\textbf{约旦标准型}——对角线为特征值、紧邻上方有 $1$ 的“约旦块”。设 $\lambda_1$ 为 $q$ 重根，对应链 $\mathbf{p}_1,\dots,\mathbf{p}_q$ 由\emph{广义特征矢量}方程递推求得\cite{liubao2006modern}：
\begin{equation}
  (\lambda_1\mathbf{I}-\mathbf{A})\mathbf{p}_1=\mathbf{0},\quad (\lambda_1\mathbf{I}-\mathbf{A})\mathbf{p}_2=-\mathbf{p}_1,\quad\dots,\quad (\lambda_1\mathbf{I}-\mathbf{A})\mathbf{p}_q=-\mathbf{p}_{q-1}.
  \label{eq:ssb-generalized-eigvec}
\end{equation}
此时 $\mathbf{T}^{-1}\mathbf{A}\mathbf{T}=\mathbf{J}$ 含一个 $q\times q$ 约旦块 $\begin{bmatrix}\lambda_1&1\\&\lambda_1&\ddots\\&&\ddots&1\\&&&\lambda_1\end{bmatrix}$。

\begin{example}[化约旦型]\label{ex:ssb-jordan}
$\mathbf{A}=\begin{bmatrix}0&1&0\\0&0&1\\2&3&0\end{bmatrix}$，$|\lambda\mathbf{I}-\mathbf{A}|=\lambda^3-3\lambda-2=0$，得 $\lambda_{1,2}=-1$（二重）、$\lambda_3=2$。$\lambda_1=-1$ 的特征矢量 $\mathbf{p}_1=[1,-1,1]^\top$；由 \cref{eq:ssb-generalized-eigvec} 求广义特征矢量 $\mathbf{p}_2=[1,0,-1]^\top$；$\lambda_3=2$ 的 $\mathbf{p}_3=[1,2,4]^\top$。则 $\mathbf{T}=[\mathbf{p}_1\,\mathbf{p}_2\,\mathbf{p}_3]$ 给出
$\mathbf{J}=\begin{bmatrix}-1&1&0\\0&-1&0\\0&0&2\end{bmatrix}$。
\end{example}

\begin{pitfall}[概念误区：约旦型“$\mathbf{J}=\mathbf{A}$”]
源教材 OCR 多处把变换后矩阵误写成 $\mathbf{J}=\mathbf{A}$ 或 $\mathbf{A}=\mathbf{T}^{-1}\mathbf{A}\mathbf{T}$\cite{liubao2006modern}——这是错的（除非 $\mathbf{A}$ 与 $\mathbf{T}$ 可交换）。变换后的矩阵是\emph{新坐标下的}系统矩阵 $\mathbf{J}=\mathbf{T}^{-1}\mathbf{A}\mathbf{T}$，记号上应与 $\mathbf{A}$ \emph{区分}（本书已统一更正）。概念上：相似变换给出同一线性映射在不同基下的矩阵，它们\emph{相似但不相等}。
\end{pitfall}

\begin{insight}
约旦型的物理意义是“\emph{解耦成模态}”。对角型（互异根）下系统彻底解耦为 $n$ 个独立的一阶模态 $\dot z_i=\lambda_i z_i$，每个按 $\mathrm{e}^{\lambda_i t}$ 独立演化——这就是为什么 $\mathrm{e}^{\boldsymbol{\Lambda}t}=\mathrm{diag}(\mathrm{e}^{\lambda_1 t},\dots)$ 平凡可写（\cref{sec:ssb-expm}）。约旦块则对应\emph{无法完全解耦}的重模态，演化里会冒出 $t\mathrm{e}^{\lambda t}$ 这样的多项式$\times$指数项（块越大、多项式次数越高）。这一“模态视角”将贯穿稳定性（看 $\mathrm{Re}\,\lambda_i$）、能控能观（看模态是否被输入激发/被输出观测，\cref{ch:controllability-obsv} 的“孤立方块”直觉）。
\end{insight}

\begin{practice}
\begin{enumerate}
  \item 验证 \cref{ex:ssb-diag} 的三个特征矢量，并算出 $\mathbf{T}^{-1}$、确认 $\mathbf{T}^{-1}\mathbf{A}\mathbf{T}=\mathrm{diag}(-1,-2,-3)$。
  \item 对 \cref{ex:ssb-jordan}，验证 $\mathbf{p}_2$ 满足广义特征矢量方程 $(\lambda_1\mathbf{I}-\mathbf{A})\mathbf{p}_2=-\mathbf{p}_1$。
  \item （友矩阵的对角化）当 $\mathbf{A}$ 是友矩阵 \cref{eq:ssb-companion} 且特征值互异时，证明对角化变换阵是范德蒙德矩阵 $\mathbf{T}=[\,\lambda_j^{\,i-1}\,]$（$i$ 行 $j$ 列）。
\end{enumerate}
\end{practice}

\subsection{从状态空间求传递函数（变换不变量的验证）}
\label{sec:ssb-tf}

由 \cref{eq:ssb-mimo} 零初值取拉氏变换：$(s\mathbf{I}-\mathbf{A})\mathbf{X}(s)=\mathbf{B}\mathbf{U}(s)$，$\mathbf{Y}(s)=\mathbf{C}\mathbf{X}(s)+\mathbf{D}\mathbf{U}(s)$，消去 $\mathbf{X}(s)$ 得\textbf{传递函数矩阵}
\begin{equation}
  \mathbf{W}(s)=\mathbf{C}(s\mathbf{I}-\mathbf{A})^{-1}\mathbf{B}+\mathbf{D}.
  \label{eq:ssb-tf}
\end{equation}
\begin{note}\textbf{OCR 勘误.}\;源 OCR 把 $\mathbf{Y}(s)=\mathbf{C}\mathbf{X}(s)+\mathbf{D}\mathbf{U}(s)$ 误成 $\mathbf{C}(s\mathbf{I}-\mathbf{A})^{-1}+\mathbf{D}\mathbf{U}(s)$（丢 $\mathbf{B}/\mathbf{X}(s)$）；本书据勘误恢复\cite{liubao2006modern}。\end{note}
把 \cref{eq:ssb-similarity} 代入 \cref{eq:ssb-tf}，利用 $(s\mathbf{I}-\mathbf{T}^{-1}\mathbf{A}\mathbf{T})^{-1}=\mathbf{T}^{-1}(s\mathbf{I}-\mathbf{A})^{-1}\mathbf{T}$，立刻得 $\mathbf{C}\mathbf{T}\cdot\mathbf{T}^{-1}(s\mathbf{I}-\mathbf{A})^{-1}\mathbf{T}\cdot\mathbf{T}^{-1}\mathbf{B}+\mathbf{D}=\mathbf{W}(s)$——传递函数确是变换不变量，呼应 \cref{thm:ssb-invariant}。

\paragraph{组合系统（辅源补充）。}
工程系统常由子系统并联、串联、反馈而成。设两子系统 $\Sigma_i:(\mathbf{A}_i,\mathbf{B}_i,\mathbf{C}_i,\mathbf{D}_i)$、传函 $\mathbf{W}_i(s)=\mathbf{C}_i(s\mathbf{I}-\mathbf{A}_i)^{-1}\mathbf{B}_i+\mathbf{D}_i$（$i=1,2$），从它们直接拼出整体状态空间，是建模的实用补丁\cite{liubao2006modern,zhang2017modern}。三种连接的差别全在“\emph{谁的输出喂给谁的输入}”——这条信号流向决定耦合块往整体 $\mathbf{A}$ 的哪个非对角位置填，于是状态拼接顺序就翻译成确定的块结构。

\noindent \textbf{并联}（$\mathbf{u}=\mathbf{u}_1=\mathbf{u}_2$、$\mathbf{y}=\mathbf{y}_1+\mathbf{y}_2$）：两子系统共享输入、输出相加，状态互不交叉，故系统矩阵\emph{块对角}、状态直接堆叠：
\begin{equation}
  \dot{\mathbf{x}}=\begin{bmatrix}\mathbf{A}_1&\mathbf{0}\\\mathbf{0}&\mathbf{A}_2\end{bmatrix}\mathbf{x}+\begin{bmatrix}\mathbf{B}_1\\\mathbf{B}_2\end{bmatrix}\mathbf{u},\quad
  \mathbf{y}=\begin{bmatrix}\mathbf{C}_1&\mathbf{C}_2\end{bmatrix}\mathbf{x}+(\mathbf{D}_1+\mathbf{D}_2)\mathbf{u},
  \label{eq:ssb-parallel}
\end{equation}
对应 $\mathbf{W}(s)=\mathbf{W}_1(s)+\mathbf{W}_2(s)$——传函相加。

\noindent \textbf{串联}（前级输出接后级输入：$\mathbf{u}=\mathbf{u}_1$、$\mathbf{u}_2=\mathbf{y}_1$、$\mathbf{y}=\mathbf{y}_2$，要求 $\dim\mathbf{y}_1=\dim\mathbf{u}_2$）：把前级输出 $\mathbf{y}_1=\mathbf{C}_1\mathbf{x}_1+\mathbf{D}_1\mathbf{u}$ 代入后级的状态方程 $\dot{\mathbf{x}}_2=\mathbf{A}_2\mathbf{x}_2+\mathbf{B}_2\mathbf{u}_2$ 与输出方程，单向耦合项 $\mathbf{B}_2\mathbf{C}_1$ 落到\emph{左下块}：
\begin{equation}
  \dot{\mathbf{x}}=\begin{bmatrix}\mathbf{A}_1&\mathbf{0}\\\mathbf{B}_2\mathbf{C}_1&\mathbf{A}_2\end{bmatrix}\mathbf{x}+\begin{bmatrix}\mathbf{B}_1\\\mathbf{B}_2\mathbf{D}_1\end{bmatrix}\mathbf{u},\quad
  \mathbf{y}=\begin{bmatrix}\mathbf{D}_2\mathbf{C}_1&\mathbf{C}_2\end{bmatrix}\mathbf{x}+\mathbf{D}_2\mathbf{D}_1\mathbf{u},
  \label{eq:ssb-conn-series}
\end{equation}
对应 $\mathbf{W}(s)=\mathbf{W}_2(s)\mathbf{W}_1(s)$（次序不可颠倒——信号\emph{先过 $\Sigma_1$ 再过 $\Sigma_2$}，故 $\mathbf{W}_2$ 在左）。

\noindent \textbf{反馈}（负反馈，前向 $\Sigma_1$、反馈 $\Sigma_2$：$\mathbf{u}_1=\mathbf{u}-\mathbf{y}_2$、$\mathbf{y}=\mathbf{y}_1=\mathbf{u}_2$；为简洁设 $\mathbf{D}_1=\mathbf{D}_2=\mathbf{0}$）：回送量 $\mathbf{y}_2=\mathbf{C}_2\mathbf{x}_2$ 经 $\mathbf{B}_1$ 反相进入前级、前级输出 $\mathbf{y}_1=\mathbf{C}_1\mathbf{x}_1$ 经 $\mathbf{B}_2$ 正向驱动反馈级，得\emph{两个}非对角耦合块：
\begin{equation}
  \dot{\mathbf{x}}=\begin{bmatrix}\mathbf{A}_1&-\mathbf{B}_1\mathbf{C}_2\\\mathbf{B}_2\mathbf{C}_1&\mathbf{A}_2\end{bmatrix}\mathbf{x}+\begin{bmatrix}\mathbf{B}_1\\\mathbf{0}\end{bmatrix}\mathbf{u},\quad
  \mathbf{y}=\begin{bmatrix}\mathbf{C}_1&\mathbf{0}\end{bmatrix}\mathbf{x}.
  \label{eq:ssb-conn-feedback}
\end{equation}
其传函须解分块求逆（由 $\mathbf{y}_1=\mathbf{W}_1(\mathbf{u}-\mathbf{W}_2\mathbf{y}_1)$ 解出）；在 $\det[\mathbf{I}+\mathbf{W}_1(s)\mathbf{W}_2(s)]\neq0$（等价 $\det[\mathbf{I}+\mathbf{W}_2(s)\mathbf{W}_1(s)]\neq0$）时，得闭环传函的两个等价写法
\begin{equation}
  \mathbf{W}(s)=\mathbf{W}_1(s)\big[\mathbf{I}+\mathbf{W}_2(s)\mathbf{W}_1(s)\big]^{-1}=\big[\mathbf{I}+\mathbf{W}_1(s)\mathbf{W}_2(s)\big]^{-1}\mathbf{W}_1(s).
  \label{eq:ssb-conn-feedback-tf}
\end{equation}
两式互等是矩阵 push-through 恒等式 $\mathbf{W}_1(\mathbf{I}+\mathbf{W}_2\mathbf{W}_1)^{-1}=(\mathbf{I}+\mathbf{W}_1\mathbf{W}_2)^{-1}\mathbf{W}_1$ 的直接结果。\emph{常数反馈} $\mathbf{y}_2=\mathbf{H}\mathbf{y}$ 是其特例：闭环 $\mathbf{A}\mapsto\mathbf{A}-\mathbf{B}\mathbf{H}\mathbf{C}$，闭环传函 $\mathbf{W}(s)=[\mathbf{I}+\mathbf{W}_0(s)\mathbf{H}]^{-1}\mathbf{W}_0(s)=\mathbf{W}_0(s)[\mathbf{I}+\mathbf{H}\mathbf{W}_0(s)]^{-1}$，其中前向通道 $\mathbf{W}_0(s)=\mathbf{C}(s\mathbf{I}-\mathbf{A})^{-1}\mathbf{B}$。

\begin{example}[2$\times$2 反馈连接的闭环传函]\label{ex:ssb-conn-feedback}
双输入双输出系统，前向通道与单位反馈 $\mathbf{H}=\mathbf{I}$\cite{zhang2017modern}：
\[
  \mathbf{W}_0(s)=\begin{bmatrix}\dfrac{1}{2s+1}&0\\[6pt]1&\dfrac{1}{s+1}\end{bmatrix},\qquad
  \mathbf{W}(s)=\big[\mathbf{I}+\mathbf{W}_0(s)\mathbf{H}\big]^{-1}\mathbf{W}_0(s).
\]
按上式，$\mathbf{I}+\mathbf{W}_0=\begin{bmatrix}\frac{2s+2}{2s+1}&0\\1&\frac{s+2}{s+1}\end{bmatrix}$，其行列式为 $\frac{2(s+2)}{2s+1}$，求逆后左乘 $\mathbf{W}_0$ 得
\[
  \mathbf{W}(s)=\begin{bmatrix}\dfrac{1}{2(s+1)}&0\\[6pt]\dfrac{2s+1}{2(s+2)}&\dfrac{1}{s+2}\end{bmatrix}.
\]
注意闭环把第一通道极点 $s=-\tfrac12$ 推到 $s=-1$、把第二通道极点 $s=-1$ 推到 $s=-2$——反馈\emph{移动了极点}，这正是 \cref{ch:linear-synthesis} 状态反馈配置极点在传函侧的影像。
\end{example}

\begin{insight}
三种连接的“块结构记法”一图记牢：\emph{并联}只填对角（子系统互不耦合）、\emph{串联}填一个非对角块 $\mathbf{B}_2\mathbf{C}_1$（单向：前级喂后级）、\emph{反馈}填两个非对角块 $\mathbf{B}_2\mathbf{C}_1$ 与 $-\mathbf{B}_1\mathbf{C}_2$（双向，一正一负，负号来自负反馈回送）。所以“拼接顺序决定块结构”不是口号——它由信号流向逐条翻译成 $\mathbf{A}$ 的非对角元。这一模块化建模是大型机器人系统搭模型的常规手段：多关节传动链 $=$ 串联、多传感器融合 $=$ 并联、带控制器的对象 $=$ 反馈；只要子系统 $(\mathbf{A}_i,\mathbf{B}_i,\mathbf{C}_i,\mathbf{D}_i)$ 已知，整机模型即可按上述三条规则机械地拼出。
\end{insight}

% ════════════════════════════════════════════════════════════════
\section{状态方程的解与矩阵指数}
\label{sec:ssb-expm}

建模告一段落，现在求解：给定 $\dot{\mathbf{x}}=\mathbf{A}\mathbf{x}+\mathbf{B}\mathbf{u}$ 与初值，状态怎么随时间演化？答案的核心是\textbf{矩阵指数} $\mathrm{e}^{\mathbf{A}t}$，它把标量解 $x(t)=\mathrm{e}^{at}x_0$ 推广到向量情形。本节先把 $\mathrm{e}^{\mathbf{A}t}$ 从直觉到定义、性质讲透，下一节给四种算法。

\subsection{齐次解与矩阵指数的来由}
\label{sec:ssb-homog}

\paragraph{从标量到矩阵：幂级数硬推。}
考虑自由运动（输入为零）$\dot{\mathbf{x}}=\mathbf{A}\mathbf{x},\ \mathbf{x}(0)=\mathbf{x}_0$。仿标量做法，设解为幂级数 $\mathbf{x}(t)=\mathbf{b}_0+\mathbf{b}_1 t+\mathbf{b}_2 t^2+\dots$，代入方程逐次匹配 $t$ 同次幂系数：
\[
  \mathbf{b}_1=\mathbf{A}\mathbf{b}_0,\quad \mathbf{b}_2=\tfrac{1}{2!}\mathbf{A}^2\mathbf{b}_0,\quad\dots,\quad \mathbf{b}_k=\tfrac{1}{k!}\mathbf{A}^k\mathbf{b}_0,
\]
而 $t=0$ 给 $\mathbf{b}_0=\mathbf{x}_0$。代回得
\begin{equation}
  \mathbf{x}(t)=\Big(\mathbf{I}+\mathbf{A}t+\tfrac{1}{2!}\mathbf{A}^2 t^2+\dots+\tfrac{1}{k!}\mathbf{A}^k t^k+\dots\Big)\mathbf{x}_0.
  \label{eq:ssb-series-sol}
\end{equation}
括号内是一个 $n\times n$ 矩阵幂级数，记为\textbf{矩阵指数}
\begin{equation}
  \boxed{\;\mathrm{e}^{\mathbf{A}t}:=\sum_{k=0}^{\infty}\frac{1}{k!}\mathbf{A}^k t^k=\mathbf{I}+\mathbf{A}t+\tfrac{1}{2!}\mathbf{A}^2t^2+\cdots\;}
  \label{eq:ssb-expm-def}
\end{equation}
于是 $\mathbf{x}(t)=\mathrm{e}^{\mathbf{A}t}\mathbf{x}_0$（初值在 $t_0$ 时为 $\mathbf{x}(t)=\mathrm{e}^{\mathbf{A}(t-t_0)}\mathbf{x}_0$）。这条幂级数推导“天然”造出了 $\mathrm{e}^{\mathbf{A}t}$——它不是符号约定，而是解的必然形式。

\begin{definition}[状态转移矩阵]\label{def:ssb-stm}
矩阵 $\boldsymbol{\Phi}(t):=\mathrm{e}^{\mathbf{A}t}$ 称\textbf{状态转移矩阵}：它把初始状态 $\mathbf{x}(0)$ 映为 $t$ 时刻状态 $\mathbf{x}(t)=\boldsymbol{\Phi}(t)\mathbf{x}(0)$，即“在状态空间里推着状态随时间转移”。它不是常矩阵——元素是 $t$ 的函数。
\end{definition}

\paragraph{几何直觉：状态轨线与组合性质。}
$\mathrm{e}^{\mathbf{A}t}$ 的几何意义最直观（图 \ref{fig:ssb-phaseflow}）：从 $\mathbf{x}(0)$ 出发，$\boldsymbol{\Phi}(t)$ 把它沿一条轨线推到 $\mathbf{x}(t)$。先转移到 $t_1$ 再转移到 $t_2$，等于直接转移到 $t_2$：$\boldsymbol{\Phi}(t_2-t_1)\boldsymbol{\Phi}(t_1)=\boldsymbol{\Phi}(t_2)$。这就是\emph{组合性质}——状态方程的解可在时间上任意分段求取，免去经典理论里处理初始条件的麻烦。

\begin{figure}[ht]
\centering
\begin{tikzpicture}[scale=1.0,>=Stealth, every node/.style={font=\small}]
% phase-plane state trajectory; explicit absolute coords, no calc-in-foreach
\draw[->] (-0.3,0)--(4.6,0) node[right]{$x_1$};
\draw[->] (0,-0.3)--(0,3.4) node[above]{$x_2$};
\draw[thick, main] (3.8,3.0) .. controls (2.6,2.4) and (1.7,1.4) .. (1.3,0.5);
\fill (3.8,3.0) circle (1.4pt) node[above right] {$\mathbf{x}(0)$};
\fill (2.55,1.95) circle (1.4pt) node[right] {$\mathbf{x}(t_1)$};
\fill (1.3,0.5) circle (1.4pt) node[below left] {$\mathbf{x}(t_2)$};
\draw[->, second, thick] (3.8,3.0) .. controls (3.3,2.7) and (2.85,2.3) .. (2.55,1.95);
\draw[->, second, thick] (2.55,1.95) .. controls (2.1,1.45) and (1.65,0.95) .. (1.3,0.5);
\node[second] at (3.55,2.35) {$\boldsymbol{\Phi}(t_1)$};
\node[second] at (1.55,1.45) {$\boldsymbol{\Phi}(t_2-t_1)$};
\end{tikzpicture}
\caption{二维状态轨线与组合性质：$\boldsymbol{\Phi}(t_1)$ 把 $\mathbf{x}(0)$ 推到 $\mathbf{x}(t_1)$，再由 $\boldsymbol{\Phi}(t_2-t_1)$ 推到 $\mathbf{x}(t_2)$，合成等于 $\boldsymbol{\Phi}(t_2)$。（仿 \cite{liubao2006modern} 重绘）}
\label{fig:ssb-phaseflow}
\end{figure}

\begin{proposition}[状态转移矩阵的基本性质]\label{prop:ssb-stm}
由定义 \cref{eq:ssb-expm-def} 可证（$\mathbf{A}$ 为定常 $n\times n$）：
\begin{align}
  &\text{(组合)}\ \boldsymbol{\Phi}(t)\boldsymbol{\Phi}(\tau)=\boldsymbol{\Phi}(t+\tau),\quad \mathrm{e}^{\mathbf{A}t}\mathrm{e}^{\mathbf{A}\tau}=\mathrm{e}^{\mathbf{A}(t+\tau)};\label{eq:ssb-stm-compose}\\
  &\text{(初值)}\ \boldsymbol{\Phi}(0)=\mathbf{I};\qquad \text{(求逆)}\ \boldsymbol{\Phi}^{-1}(t)=\boldsymbol{\Phi}(-t)=\mathrm{e}^{-\mathbf{A}t};\label{eq:ssb-stm-inv}\\
  &\text{(微分)}\ \dot{\boldsymbol{\Phi}}(t)=\mathbf{A}\boldsymbol{\Phi}(t)=\boldsymbol{\Phi}(t)\mathbf{A};\label{eq:ssb-stm-deriv}\\
  &\text{(可乘性，受限)}\ \mathrm{e}^{\mathbf{A}t}\mathrm{e}^{\mathbf{B}t}=\mathrm{e}^{(\mathbf{A}+\mathbf{B})t}\ \Longleftrightarrow\ \mathbf{A}\mathbf{B}=\mathbf{B}\mathbf{A}.\label{eq:ssb-stm-commute}
\end{align}
\end{proposition}
\begin{proof}[证（要点）]
组合性、初值、微分均由 \cref{eq:ssb-expm-def} 逐项求导/相乘并比较级数得到（如 $\dot{\boldsymbol{\Phi}}=\mathbf{A}+\mathbf{A}^2 t+\dots=\mathbf{A}(\mathbf{I}+\mathbf{A}t+\dots)=\mathbf{A}\mathrm{e}^{\mathbf{A}t}$，又 $=\mathrm{e}^{\mathbf{A}t}\mathbf{A}$ 因 $\mathbf{A}$ 与自身幂可交换）。求逆由组合性取 $\tau=-t$ 加初值得。可乘性：把 $\mathrm{e}^{\mathbf{A}t}\mathrm{e}^{\mathbf{B}t}$ 与 $\mathrm{e}^{(\mathbf{A}+\mathbf{B})t}$ 展到 $t^2$ 比较，$t^2$ 项分别含 $\tfrac12(\mathbf{A}^2+2\mathbf{A}\mathbf{B}+\mathbf{B}^2)$ 与 $\tfrac12(\mathbf{A}+\mathbf{B})^2=\tfrac12(\mathbf{A}^2+\mathbf{A}\mathbf{B}+\mathbf{B}\mathbf{A}+\mathbf{B}^2)$，相等当且仅当 $\mathbf{A}\mathbf{B}=\mathbf{B}\mathbf{A}$。
\end{proof}

\begin{pitfall}[思维陷阱：把 $\mathrm{e}^{\mathbf{A}t}$ 当标量指数无脑用]
标量恒等式 $\mathrm{e}^{a}\mathrm{e}^{b}=\mathrm{e}^{a+b}$ 对矩阵\emph{一般不成立}——\cref{eq:ssb-stm-commute} 要求 $\mathbf{A},\mathbf{B}$ 可交换。矩阵不可交换是常态，所以见到 $\mathrm{e}^{\mathbf{A}t}\mathrm{e}^{\mathbf{B}t}$ 切勿擅自合并为 $\mathrm{e}^{(\mathbf{A}+\mathbf{B})t}$。这一“可交换性壁垒”在时变系统（\cref{sec:ssb-ltv}）会再次现身，且更要命：那里连 $\mathrm{e}^{\int\mathbf{A}(\tau)\mathrm{d}\tau}$ 都未必是解。同理：李群（\cref{ch:lie}）的 BCH 公式正是处理 $\mathrm{e}^{\mathbf{A}}\mathrm{e}^{\mathbf{B}}$ 的“修正项”，本质同源。
\end{pitfall}

\paragraph{几个可直接写出的特例。}
\begin{itemize}
  \item 对角阵 $\boldsymbol{\Lambda}=\mathrm{diag}(\lambda_i)$：$\mathrm{e}^{\boldsymbol{\Lambda}t}=\mathrm{diag}(\mathrm{e}^{\lambda_i t})$；
  \item 可对角化 $\mathbf{T}^{-1}\mathbf{A}\mathbf{T}=\boldsymbol{\Lambda}$：$\mathrm{e}^{\mathbf{A}t}=\mathbf{T}\,\mathrm{e}^{\boldsymbol{\Lambda}t}\,\mathbf{T}^{-1}$；
  \item 约旦块 $\mathbf{J}$（特征值 $\lambda$、阶 $n$）：$\mathrm{e}^{\mathbf{J}t}=\mathrm{e}^{\lambda t}\!\begin{bmatrix}1&t&\tfrac{t^2}{2!}&\cdots&\tfrac{t^{n-1}}{(n-1)!}\\0&1&t&\cdots&\tfrac{t^{n-2}}{(n-2)!}\\\vdots&&\ddots&&\vdots\\0&0&\cdots&1&t\\0&0&\cdots&0&1\end{bmatrix}$（多项式$\times$指数，印证约旦块“耦合模态”的物理意义）；
  \item 复模态块 $\mathbf{A}=\begin{bmatrix}\sigma&\omega\\-\omega&\sigma\end{bmatrix}$：$\mathrm{e}^{\mathbf{A}t}=\mathrm{e}^{\sigma t}\begin{bmatrix}\cos\omega t&\sin\omega t\\-\sin\omega t&\cos\omega t\end{bmatrix}$（衰减/增长的旋转）。
\end{itemize}

% ════════════════════════════════════════════════════════════════
\section{计算 \texorpdfstring{$\mathrm{e}^{\mathbf{A}t}$}{e\^{}\{At\}} 的四种算法}
\label{sec:ssb-four-algos}

给定 $\mathbf{A}$，如何把 $\mathrm{e}^{\mathbf{A}t}$ 算出来？四种主流算法各擅胜场：定义级数（编程友好、得不到闭式）、约旦型（结构清晰、需求 $\mathbf{T}$）、拉氏反变换（系统化、SISO 尤便）、Cayley--Hamilton（降幂、解析）。本节以同一个 $\mathbf{A}=\begin{bmatrix}0&1\\-2&-3\end{bmatrix}$ 贯穿四法对照。

\subsection{算法一：定义级数（数值）}
直接截断 \cref{eq:ssb-expm-def}：
\[
  \mathrm{e}^{\mathbf{A}t}=\mathbf{I}+\mathbf{A}t+\tfrac{1}{2!}\mathbf{A}^2t^2+\tfrac{1}{3!}\mathbf{A}^3t^3+\cdots
  =\begin{bmatrix}1-t^2+t^3+\cdots & t-\tfrac{3}{2}t^2+\tfrac{7}{6}t^3+\cdots\\[2pt] -2t+3t^2-\tfrac{7}{3}t^3+\cdots & 1-3t+\tfrac{7}{2}t^2-\tfrac{5}{2}t^3+\cdots\end{bmatrix}.
\]
\begin{note}\textbf{OCR 勘误.}\;源 OCR 把 $(1,2)$ 元 $t^3$ 系数的符号识反（印成 $-\tfrac76 t^3$）；按 $\mathbf{A}^3$ 直接展开该系数为 $+\tfrac76$，本书已更正\cite{liubao2006modern}。\end{note}
优点：步骤简单、易编程，适合计算机。缺点：一般得不到\emph{解析}闭式，且高维收敛慢。实践中数值算 $\mathrm{e}^{\mathbf{A}t}$ 多用 scaling-and-squaring 或 Padé 近似，但原理同源。

\subsection{算法二：化约旦标准型}
若 $\mathbf{A}$ 互异根，$\boldsymbol{\Lambda}=\mathbf{T}^{-1}\mathbf{A}\mathbf{T}$，则 $\mathrm{e}^{\mathbf{A}t}=\mathbf{T}\,\mathrm{e}^{\boldsymbol{\Lambda}t}\,\mathbf{T}^{-1}$。
\begin{note}\textbf{OCR 勘误.}\;源 OCR 把中间因子误写成 $\mathrm{e}^{\mathbf{A}t}=\mathbf{T}\mathrm{e}^{\mathbf{A}t}\mathbf{T}^{-1}$（自相矛盾）；正确为 $\mathbf{T}\mathrm{e}^{\boldsymbol{\Lambda}t}\mathbf{T}^{-1}$\cite{liubao2006modern}。\end{note}
对 $\mathbf{A}=\begin{bmatrix}0&1\\-2&-3\end{bmatrix}$：$\lambda_1=-1,\lambda_2=-2$，取 $\mathbf{T}=\begin{bmatrix}2&1\\-2&-2\end{bmatrix}$、$\mathbf{T}^{-1}=\begin{bmatrix}1&\tfrac12\\-1&-1\end{bmatrix}$，得
\begin{equation}
  \mathrm{e}^{\mathbf{A}t}=\mathbf{T}\begin{bmatrix}\mathrm{e}^{-t}&0\\0&\mathrm{e}^{-2t}\end{bmatrix}\mathbf{T}^{-1}
  =\begin{bmatrix}2\mathrm{e}^{-t}-\mathrm{e}^{-2t}&\mathrm{e}^{-t}-\mathrm{e}^{-2t}\\-2\mathrm{e}^{-t}+2\mathrm{e}^{-2t}&-\mathrm{e}^{-t}+2\mathrm{e}^{-2t}\end{bmatrix}.
  \label{eq:ssb-expm-result}
\end{equation}
重根时改用 $\mathrm{e}^{\mathbf{J}t}$（含 $t\mathrm{e}^{\lambda t}$ 项）。优点：结构最清晰、模态一目了然；缺点：要先求 $\mathbf{T},\mathbf{T}^{-1}$。

\subsection{算法三：拉氏反变换}
\begin{theorem}[拉氏反变换法]\label{thm:ssb-laplace}
$\displaystyle \mathrm{e}^{\mathbf{A}t}=\mathcal{L}^{-1}\!\big[(s\mathbf{I}-\mathbf{A})^{-1}\big]$.
\end{theorem}
\begin{proof}
对 $\dot{\mathbf{x}}=\mathbf{A}\mathbf{x},\ \mathbf{x}(0)=\mathbf{x}_0$ 取拉氏变换：$s\mathbf{X}(s)-\mathbf{x}_0=\mathbf{A}\mathbf{X}(s)$，即 $\mathbf{X}(s)=(s\mathbf{I}-\mathbf{A})^{-1}\mathbf{x}_0$。反变换得 $\mathbf{x}(t)=\mathcal{L}^{-1}[(s\mathbf{I}-\mathbf{A})^{-1}]\mathbf{x}_0$，与 $\mathbf{x}(t)=\mathrm{e}^{\mathbf{A}t}\mathbf{x}_0$ 比较即证。
\end{proof}
同例：$(s\mathbf{I}-\mathbf{A})^{-1}=\dfrac{1}{(s+1)(s+2)}\begin{bmatrix}s+3&1\\-2&s\end{bmatrix}$，部分分式后逐元反变换得 \cref{eq:ssb-expm-result}。优点：系统化、SISO 尤便、顺带得到极点；缺点：高维求逆与部分分式繁琐。

\subsection{算法四：Cayley--Hamilton 降幂}
\begin{theorem}[Cayley--Hamilton 化有限项]\label{thm:ssb-cayley}
方阵 $\mathbf{A}$ 满足自身特征方程 $\mathbf{A}^n+a_{n-1}\mathbf{A}^{n-1}+\dots+a_0\mathbf{I}=\mathbf{0}$，故 $\mathbf{A}^n$（及更高幂）可由 $\{\mathbf{I},\mathbf{A},\dots,\mathbf{A}^{n-1}\}$ 线性表出。于是 $\mathrm{e}^{\mathbf{A}t}$ 中所有 $n$ 次及以上幂可消去，化为有限项：
\begin{equation}
  \mathrm{e}^{\mathbf{A}t}=\alpha_0(t)\mathbf{I}+\alpha_1(t)\mathbf{A}+\dots+\alpha_{n-1}(t)\mathbf{A}^{n-1}.
  \label{eq:ssb-cayley}
\end{equation}
\end{theorem}
系数 $\alpha_i(t)$ 由“$\lambda$ 也满足同一关系”确定：特征值互异时解范德蒙德方程组
\begin{equation}
  \begin{bmatrix}\alpha_0\\\vdots\\\alpha_{n-1}\end{bmatrix}
  =\begin{bmatrix}1&\lambda_1&\cdots&\lambda_1^{n-1}\\\vdots&\vdots&&\vdots\\1&\lambda_n&\cdots&\lambda_n^{n-1}\end{bmatrix}^{-1}
  \begin{bmatrix}\mathrm{e}^{\lambda_1 t}\\\vdots\\\mathrm{e}^{\lambda_n t}\end{bmatrix},
  \label{eq:ssb-cayley-vander}
\end{equation}
重根时对 $\lambda$ 求导补足方程。同例（$\lambda_1=-1,\lambda_2=-2$）：$\alpha_0=2\mathrm{e}^{-t}-\mathrm{e}^{-2t},\ \alpha_1=\mathrm{e}^{-t}-\mathrm{e}^{-2t}$，代入 \cref{eq:ssb-cayley} 仍得 \cref{eq:ssb-expm-result}。
\begin{note}\textbf{OCR 勘误.}\;源 OCR 在 Cayley--Hamilton 式把 $a_0\mathbf{I}$ 误加撇号（$a_0'\mathbf{I}$），并把 $\alpha_0(t)$ 级数 $t^3$ 符号识错；本书据勘误更正\cite{liubao2006modern}。\end{note}
优点：得解析式、只需 $\mathbf{A}$ 的幂与特征值；缺点：高维范德蒙德求逆病态。

\begin{insight}
四法殊途同归（都得 \cref{eq:ssb-expm-result}），但“为什么都对”值得品味：它们本质都在做同一件事——把 $\mathrm{e}^{\mathbf{A}t}$ 表达成 $\mathbf{A}$ 的特征值上的标量函数 $\mathrm{e}^{\lambda_i t}$ 的某种组合。约旦型直接在模态坐标做；Cayley--Hamilton 用“$\mathbf{A}$ 与 $\lambda$ 满足同一最小多项式”把矩阵函数化为标量函数；拉氏法则在频域做部分分式（极点 $=\lambda_i$）。\emph{矩阵函数 $=$ 谱上的标量函数}是统一主线。工程选择：低维手算用拉氏或 Cayley--Hamilton（要解析式）；高维上机用数值级数/Padé；要看清模态结构用约旦型。
\end{insight}

\begin{practice}
\begin{enumerate}
  \item 用四种算法分别求 $\mathbf{A}=\begin{bmatrix}0&1\\0&-2\end{bmatrix}$ 的 $\mathrm{e}^{\mathbf{A}t}$，核对一致（提示：$\lambda=0,-2$）。
  \item 对含二重根的 $\mathbf{A}=\begin{bmatrix}0&1&0\\0&0&1\\2&-5&4\end{bmatrix}$（$\lambda_{1,2}=1,\lambda_3=2$），用 \cref{eq:ssb-cayley-vander}（重根补求导方程）求 $\alpha_i(t)$。
  \item 证明 \cref{prop:ssb-stm} 的 \cref{eq:ssb-stm-commute} 在 $\mathbf{A}\mathbf{B}\neq\mathbf{B}\mathbf{A}$ 时给出反例（取两个不可交换的 $2\times2$ 矩阵验证两边到 $t^2$ 项不等）。
\end{enumerate}
\end{practice}

% ════════════════════════════════════════════════════════════════
\section{非齐次方程的解：自由运动加强迫运动}
\label{sec:ssb-nonhomog}

加入输入 $\mathbf{u}(t)$ 后，状态方程 $\dot{\mathbf{x}}=\mathbf{A}\mathbf{x}+\mathbf{B}\mathbf{u}$ 是非齐次的。其解由“初值引起的自由运动”加“输入激励引起的强迫运动”两部分组成。

\begin{theorem}[线性定常非齐次方程的解]\label{thm:ssb-nonhomog}
$\dot{\mathbf{x}}=\mathbf{A}\mathbf{x}+\mathbf{B}\mathbf{u}$、初值 $\mathbf{x}(t_0)$，解为
\begin{equation}
  \mathbf{x}(t)=\underbrace{\mathrm{e}^{\mathbf{A}(t-t_0)}\mathbf{x}(t_0)}_{\text{自由运动}}+\underbrace{\int_{t_0}^{t}\mathrm{e}^{\mathbf{A}(t-\tau)}\mathbf{B}\mathbf{u}(\tau)\,\mathrm{d}\tau}_{\text{强迫运动（卷积）}}.
  \label{eq:ssb-nonhomog}
\end{equation}
\end{theorem}
\begin{proof}
把方程写成 $\dot{\mathbf{x}}-\mathbf{A}\mathbf{x}=\mathbf{B}\mathbf{u}$，两边左乘积分因子 $\mathrm{e}^{-\mathbf{A}t}$。利用 $\tfrac{\mathrm{d}}{\mathrm{d}t}(\mathrm{e}^{-\mathbf{A}t}\mathbf{x})=\mathrm{e}^{-\mathbf{A}t}(\dot{\mathbf{x}}-\mathbf{A}\mathbf{x})$，得 $\tfrac{\mathrm{d}}{\mathrm{d}t}(\mathrm{e}^{-\mathbf{A}t}\mathbf{x})=\mathrm{e}^{-\mathbf{A}t}\mathbf{B}\mathbf{u}$。在 $[t_0,t]$ 积分：$\mathrm{e}^{-\mathbf{A}t}\mathbf{x}(t)-\mathrm{e}^{-\mathbf{A}t_0}\mathbf{x}(t_0)=\int_{t_0}^t\mathrm{e}^{-\mathbf{A}\tau}\mathbf{B}\mathbf{u}(\tau)\mathrm{d}\tau$，左乘 $\mathrm{e}^{\mathbf{A}t}$ 整理即得 \cref{eq:ssb-nonhomog}。（亦可由拉氏变换：$\mathbf{X}(s)=(s\mathbf{I}-\mathbf{A})^{-1}\mathbf{x}_0+(s\mathbf{I}-\mathbf{A})^{-1}\mathbf{B}\mathbf{U}(s)$，第二项是 $\boldsymbol{\Phi}\ast\mathbf{B}\mathbf{u}$ 的拉氏，反变换得卷积。）
\end{proof}

\begin{insight}
\cref{eq:ssb-nonhomog} 的结构是整个线性系统理论的“主旋律”：解 $=$ 零输入响应 $+$ 零状态响应。第一项只看初值、与输入无关；第二项是\emph{冲激响应 $\mathrm{e}^{\mathbf{A}t}\mathbf{B}$ 与输入的卷积}——这正是“线性时不变系统由其冲激响应完全刻画”的状态空间版本。能控性（\cref{ch:controllability-obsv}）问的就是：调度 $\mathbf{u}(\tau)$ 能否让这个卷积积分把状态送到任意目标；LQR（\cref{sec:ctrl-lqr}）则是在这个解的基础上优化 $\mathbf{u}$。所以这条公式是承上启下的枢纽。
\end{insight}

\paragraph{典型输入的简化（要点）。}
对常见输入，卷积可显式算出\cite{liubao2006modern}：冲激 $\mathbf{u}=\mathbf{K}\delta(t)$ 给 $\mathbf{x}=\mathrm{e}^{\mathbf{A}t}\mathbf{x}_0+\mathrm{e}^{\mathbf{A}t}\mathbf{B}\mathbf{K}$；阶跃 $\mathbf{u}=\mathbf{K}\cdot\mathbf{1}(t)$ 给 $\mathbf{x}=\mathrm{e}^{\mathbf{A}t}\mathbf{x}_0+\mathbf{A}^{-1}(\mathrm{e}^{\mathbf{A}t}-\mathbf{I})\mathbf{B}\mathbf{K}$（$\mathbf{A}$ 可逆时）；斜坡 $\mathbf{u}=\mathbf{K}t\cdot\mathbf{1}(t)$ 给 $\mathbf{x}=\mathrm{e}^{\mathbf{A}t}\mathbf{x}_0+[\mathbf{A}^{-2}(\mathrm{e}^{\mathbf{A}t}-\mathbf{I})-\mathbf{A}^{-1}t]\mathbf{B}\mathbf{K}$。
\begin{note}\textbf{OCR 勘误.}\;斜坡解的源 OCR 把标量因子 $\mathbf{A}^{-1}t$ 误并入上标（印成 $\mathbf{A}^{-1\,t}$）；$t$ 是乘性标量、非指数，本书已更正\cite{liubao2006modern}。\end{note}

% ════════════════════════════════════════════════════════════════
\section{线性时变系统的解（辅源补强）}
\label{sec:ssb-ltv}

机器人沿一条名义轨迹做线性化，得到的是\emph{线性时变（LTV）}系统 $\dot{\mathbf{x}}=\mathbf{A}(t)\mathbf{x}+\mathbf{B}(t)\mathbf{u}$——这是 iLQR/DDP、轨迹跟踪、时变 LQR 的核心模型。本节把定常理论推广到时变，关键认知是：定常的“$\mathrm{e}^{\mathbf{A}t}$ 闭式”一般\emph{不再可用}，须诉诸 Peano--Baker 级数或基本解阵。本节以辅源\cite{zhang2017modern}讲法为主线（它把可交换条件与 Peano--Baker 讲得更透）。

\paragraph{反面：可交换条件为何苛刻。}
标量时变 $\dot x=a(t)x$ 由分离变量得 $x(t)=\mathrm{e}^{\int_{t_0}^t a(\tau)\mathrm{d}\tau}x(t_0)$。能否直接推广为 $\mathbf{x}(t)=\mathrm{e}^{\int_{t_0}^t\mathbf{A}(\tau)\mathrm{d}\tau}\mathbf{x}(t_0)$？

\begin{theorem}[时变指数解的成立条件]\label{thm:ssb-ltv-commute}
$\mathbf{x}(t)=\mathrm{e}^{\int_{t_0}^t\mathbf{A}(\tau)\mathrm{d}\tau}\mathbf{x}(t_0)$ 是 $\dot{\mathbf{x}}=\mathbf{A}(t)\mathbf{x}$ 的解，当且仅当 $\mathbf{A}(t)$ 与 $\int_{t_0}^t\mathbf{A}(\tau)\mathrm{d}\tau$ 矩阵乘法可交换：
\begin{equation}
  \mathbf{A}(t)\!\int_{t_0}^t\!\mathbf{A}(\tau)\mathrm{d}\tau=\int_{t_0}^t\!\mathbf{A}(\tau)\mathrm{d}\tau\,\mathbf{A}(t).
  \label{eq:ssb-ltv-commute}
\end{equation}
\end{theorem}
\begin{proof}[证（要点）]
把 $\mathrm{e}^{\int\mathbf{A}}=\mathbf{I}+\int_{t_0}^t\mathbf{A}\,\mathrm{d}\tau+\tfrac{1}{2!}(\int_{t_0}^t\mathbf{A}\,\mathrm{d}\tau)^2+\dots$ 对 $t$ 求导。逐项用乘积法则，二次项导数含 $\tfrac12[\mathbf{A}(t)\!\int\mathbf{A}+\int\mathbf{A}\,\mathbf{A}(t)]$；而要满足 $\tfrac{\mathrm{d}}{\mathrm{d}t}\mathrm{e}^{\int\mathbf{A}}=\mathbf{A}(t)\mathrm{e}^{\int\mathbf{A}}$，右端二次项是 $\mathbf{A}(t)\!\int\mathbf{A}$。两者相等 $\iff$ \cref{eq:ssb-ltv-commute}。此条件\emph{一般不成立}（如 $\mathbf{A}(t_1)\mathbf{A}(t_2)\neq\mathbf{A}(t_2)\mathbf{A}(t_1)$ 时），故时变自由解通常无封闭形式。
\end{proof}

\begin{example}[可交换 vs 不可交换]\label{ex:ssb-ltv-commute}
\emph{可交换}：$\mathbf{A}(t)=\begin{bmatrix}0&\tfrac{1}{(t+1)^2}\\0&0\end{bmatrix}$，任意 $t_1,t_2$ 有 $\mathbf{A}(t_1)\mathbf{A}(t_2)=\mathbf{0}=\mathbf{A}(t_2)\mathbf{A}(t_1)$，可用 \cref{eq:ssb-ltv-commute} 的指数式，解得 $\mathbf{x}(t)=\begin{bmatrix}1&t-t_0\\0&1\end{bmatrix}\mathbf{x}(t_0)$。\emph{不可交换}：$\mathbf{A}(t)=\begin{bmatrix}0&1\\0&t\end{bmatrix}$，$\mathbf{A}(t_1)\mathbf{A}(t_2)\neq\mathbf{A}(t_2)\mathbf{A}(t_1)$，必须用下面的 Peano--Baker 级数。
\end{example}

\paragraph{理论：Peano--Baker 级数与状态转移矩阵。}
\begin{theorem}[Peano--Baker 级数]\label{thm:ssb-peano-baker}
不论可交换与否，$\dot{\mathbf{x}}=\mathbf{A}(t)\mathbf{x}$ 的解 $\mathbf{x}(t)=\boldsymbol{\Phi}(t,t_0)\mathbf{x}(t_0)$，其中状态转移矩阵
\begin{equation}
  \boldsymbol{\Phi}(t,t_0)=\mathbf{I}+\int_{t_0}^t\!\mathbf{A}(\tau_1)\mathrm{d}\tau_1+\int_{t_0}^t\!\mathbf{A}(\tau_1)\!\int_{t_0}^{\tau_1}\!\mathbf{A}(\tau_2)\,\mathrm{d}\tau_2\,\mathrm{d}\tau_1+\cdots
  \label{eq:ssb-peano-baker}
\end{equation}
（$\mathbf{A}(t)$ 元素有界时级数收敛）。它是矩阵微分方程 $\dot{\boldsymbol{\Phi}}(t,t_0)=\mathbf{A}(t)\boldsymbol{\Phi}(t,t_0),\ \boldsymbol{\Phi}(t_0,t_0)=\mathbf{I}$ 的唯一解。
\end{theorem}
\begin{proof}[证（要点）]
逐项对 $t$ 求导：第 $k$ 项的导数是把外层积分上限代入 $t$，提出公共左因子 $\mathbf{A}(t)$，恰好得到第 $k-1$ 项乘 $\mathbf{A}(t)$，故 $\dot{\boldsymbol{\Phi}}=\mathbf{A}(t)\boldsymbol{\Phi}$；$t=t_0$ 时各积分为零，$\boldsymbol{\Phi}(t_0,t_0)=\mathbf{I}$。唯一性由线性 ODE 解的唯一性得。
\end{proof}

\begin{definition}[时变状态转移矩阵的性质]\label{def:ssb-ltv-stm}
$\boldsymbol{\Phi}(t,t_0)$ 依赖\emph{两个}时刻（不仅是时间差），满足：$\boldsymbol{\Phi}(t,t)=\mathbf{I}$；$\boldsymbol{\Phi}^{-1}(t,t_0)=\boldsymbol{\Phi}(t_0,t)$；分段组合 $\boldsymbol{\Phi}(t_2,t_0)=\boldsymbol{\Phi}(t_2,t_1)\boldsymbol{\Phi}(t_1,t_0)$。若 $\boldsymbol{\Psi}(t)$ 是任一基本解阵（$n$ 个线性无关解为列），则 $\boldsymbol{\Phi}(t,t_0)=\boldsymbol{\Psi}(t)\boldsymbol{\Psi}^{-1}(t_0)$（与基本解阵选取无关）。
\end{definition}

\begin{note}
\textbf{定常是时变的特例.}\;定常时 $\mathbf{A}$ 与 $\int\mathbf{A}\,\mathrm{d}\tau=\mathbf{A}(t-t_0)$ 显然可交换，Peano--Baker 级数 \cref{eq:ssb-peano-baker} 退化为 $\sum_k\tfrac{1}{k!}\mathbf{A}^k(t-t_0)^k=\mathrm{e}^{\mathbf{A}(t-t_0)}$，且 $\boldsymbol{\Phi}(t,t_0)=\boldsymbol{\Phi}(t-t_0)$ 只依赖时间差——与前几节完全一致。这也解释了为何定常情形“好算”：可交换性使无穷级数坍缩为指数闭式。
\end{note}

\paragraph{时变非齐次解。}
有了 $\boldsymbol{\Phi}(t,t_0)$，时变非齐次方程 $\dot{\mathbf{x}}=\mathbf{A}(t)\mathbf{x}+\mathbf{B}(t)\mathbf{u}$ 的解与定常 \cref{eq:ssb-nonhomog} 形式上一致\cite{zhang2017modern}：
\begin{equation}
  \mathbf{x}(t)=\boldsymbol{\Phi}(t,t_0)\mathbf{x}_0+\int_{t_0}^t\boldsymbol{\Phi}(t,\tau)\mathbf{B}(\tau)\mathbf{u}(\tau)\,\mathrm{d}\tau.
  \label{eq:ssb-ltv-nonhomog}
\end{equation}
唯一区别是用 $\boldsymbol{\Phi}(t,\tau)$（两时刻）替换 $\mathrm{e}^{\mathbf{A}(t-\tau)}$。这种“表述上的统一”对理论分析极有利；但计算上时变 $\boldsymbol{\Phi}$ 除极简单情形外难求闭式，工程上多用数值积分。

\begin{insight}
LTV 是连接“线性定常”与“非线性”的桥。机器人本质非线性（\cref{eq:ssb-mechanical-general} 的 $\mathbf{M}(\mathbf{q})$ 依赖位形），但沿一条名义轨迹 $\bar{\mathbf{x}}(t),\bar{\mathbf{u}}(t)$ 做一阶 Taylor 展开，误差动态 $\delta\dot{\mathbf{x}}=\mathbf{A}(t)\delta\mathbf{x}+\mathbf{B}(t)\delta\mathbf{u}$ 就是 LTV，其中 $\mathbf{A}(t)=\left.\tfrac{\partial\mathbf{f}}{\partial\mathbf{x}}\right|_{\bar{\mathbf{x}},\bar{\mathbf{u}}}$、$\mathbf{B}(t)=\left.\tfrac{\partial\mathbf{f}}{\partial\mathbf{u}}\right|_{\bar{\mathbf{x}},\bar{\mathbf{u}}}$。iLQR/DDP（规划侧见 \cref{ch:control} 的 \cref{sec:ctrl-lqr}、理论深化见 A--Q）正是在每次迭代沿当前轨迹得到这样的 LTV，再对其求时变 LQR。所以本节不是“选读”——它是把线性工具用到真实机器人上的必经关口。
\end{insight}

% ════════════════════════════════════════════════════════════════
\section{连续系统的离散化}
\label{sec:ssb-discretize}

数字控制器、卡尔曼滤波、MPC 预测方程都工作在离散时间。把连续模型 $\dot{\mathbf{x}}=\mathbf{A}\mathbf{x}+\mathbf{B}\mathbf{u}$ 化为离散 $\mathbf{x}_{k+1}=\mathbf{A}_d\mathbf{x}_k+\mathbf{B}_d\mathbf{u}_k$ 是工程落地的最后一公里。关键假设是\textbf{零阶保持（ZOH）}：输入在采样间隔 $[kT,(k+1)T)$ 内保持常值 $\mathbf{u}(kT)$。

\begin{theorem}[零阶保持离散化（精确）]\label{thm:ssb-zoh}
在 ZOH 假设下，连续系统 \cref{eq:ssb-mimo} 的精确离散化为 $\mathbf{x}((k+1)T)=\mathbf{A}_d\mathbf{x}(kT)+\mathbf{B}_d\mathbf{u}(kT)$，其中
\begin{equation}
  \mathbf{A}_d=\mathrm{e}^{\mathbf{A}T},\qquad \mathbf{B}_d=\Big(\int_0^T\mathrm{e}^{\mathbf{A}t}\,\mathrm{d}t\Big)\mathbf{B},
  \label{eq:ssb-zoh}
\end{equation}
而 $\mathbf{C},\mathbf{D}$ 不变。
\end{theorem}
\begin{proof}
从连续解 \cref{eq:ssb-nonhomog} 取 $t_0=kT,\ t=(k+1)T$，并用 $\mathbf{u}(\tau)\equiv\mathbf{u}(kT)$：
\[
  \mathbf{x}((k{+}1)T)=\mathrm{e}^{\mathbf{A}T}\mathbf{x}(kT)+\Big(\int_{kT}^{(k+1)T}\mathrm{e}^{\mathbf{A}[(k+1)T-\tau]}\mathrm{d}\tau\Big)\mathbf{B}\,\mathbf{u}(kT).
\]
对积分换元 $t=(k{+}1)T-\tau$（$\mathrm{d}\tau=-\mathrm{d}t$，上下限 $\tau{:}\,kT{\to}(k{+}1)T$ 对应 $t{:}\,T{\to}0$）得 $\int_0^T\mathrm{e}^{\mathbf{A}t}\mathrm{d}t$，即 \cref{eq:ssb-zoh}。输出方程是状态/输入的线性组合，离散化不改其系数，故 $\mathbf{C},\mathbf{D}$ 不变。
\begin{note}\textbf{OCR 勘误.}\;源 OCR 在该换元式把积分变量 $\mathrm{e}^{\mathbf{A}t}$ 误印成 $\mathrm{e}^{\mathbf{A}T}$（$t\leftrightarrow T$ 混淆）；近似式中 $\mathbf{x}((k{+}1)T)$ 的离散索引 $1$ 被识成 $I$；本书均据勘误更正\cite{liubao2006modern}。\end{note}
\end{proof}

\begin{note}
\textbf{$\mathbf{B}_d$ 的可逆闭式与计算.}\;$\mathbf{A}$ 可逆时 $\int_0^T\mathrm{e}^{\mathbf{A}t}\mathrm{d}t=\mathbf{A}^{-1}(\mathrm{e}^{\mathbf{A}T}-\mathbf{I})$，故 $\mathbf{B}_d=\mathbf{A}^{-1}(\mathrm{e}^{\mathbf{A}T}-\mathbf{I})\mathbf{B}$。$\mathbf{A}$ 奇异（如含纯积分器）时用增广矩阵指数一次性算出：$\exp\!\Big(\begin{bsmallmatrix}\mathbf{A}&\mathbf{B}\\\mathbf{0}&\mathbf{0}\end{bsmallmatrix}T\Big)=\begin{bsmallmatrix}\mathbf{A}_d&\mathbf{B}_d\\\mathbf{0}&\mathbf{I}\end{bsmallmatrix}$（Van~Loan 法）——这是数值库 \texttt{c2d}/\allowbreak\texttt{cont2discrete} 的内核。
\end{note}

\begin{theorem}[一阶近似（前向欧拉）]\label{thm:ssb-euler}
采样周期 $T$ 足够小（约为系统最小时间常数的 $1/10$）时，可近似
\begin{equation}
  \mathbf{A}_d\approx\mathbf{I}+\mathbf{A}T,\qquad \mathbf{B}_d\approx\mathbf{B}T.
  \label{eq:ssb-euler}
\end{equation}
\end{theorem}
\begin{proof}
由导数定义 $\dot{\mathbf{x}}(kT)\approx\dfrac{\mathbf{x}((k{+}1)T)-\mathbf{x}(kT)}{T}$，代入 $\dot{\mathbf{x}}=\mathbf{A}\mathbf{x}+\mathbf{B}\mathbf{u}$ 整理即得；等价于把 $\mathrm{e}^{\mathbf{A}T}=\mathbf{I}+\mathbf{A}T+O(T^2)$ 截断到一阶。
\end{proof}

\begin{example}[ZOH 精确 vs 欧拉近似]\label{ex:ssb-discretize}
$\dot{\mathbf{x}}=\begin{bmatrix}0&1\\0&-2\end{bmatrix}\mathbf{x}+\begin{bmatrix}0\\1\end{bmatrix}u$。由 $\mathrm{e}^{\mathbf{A}t}=\begin{bmatrix}1&\tfrac12(1-\mathrm{e}^{-2t})\\0&\mathrm{e}^{-2t}\end{bmatrix}$ 得精确
$\mathbf{A}_d=\begin{bmatrix}1&\tfrac12(1-\mathrm{e}^{-2T})\\0&\mathrm{e}^{-2T}\end{bmatrix}$、$\mathbf{B}_d=\begin{bmatrix}\tfrac12(T+\tfrac{\mathrm{e}^{-2T}-1}{2})\\\tfrac12(1-\mathrm{e}^{-2T})\end{bmatrix}$；欧拉近似 $\mathbf{A}_d\approx\begin{bmatrix}1&T\\0&1-2T\end{bmatrix}$、$\mathbf{B}_d\approx\begin{bmatrix}0\\T\end{bmatrix}$。下表显示 $T$ 越小两者越接近、$T=0.05$ 时已几乎重合。
\begin{table}[H]\centering\small
\caption{采样周期对离散化的影响（精确 ZOH vs 一阶近似的 $\mathbf{A}_d$ 第 $(2,2)$ 元与 $\mathbf{B}_d$ 第 2 元）}
\label{tab:ssb-discretize}
\begin{tabularx}{\linewidth}{cLLLL}
\toprule
$T$ & $\mathrm{e}^{-2T}$（精确 $A_{d,22}$）& $1-2T$（近似）& $\tfrac12(1-\mathrm{e}^{-2T})$（精确 $B_{d,2}$）& $T$（近似）\\
\midrule
$1.0$  & $0.135$ & $-1.0$ & $0.432$  & $1.00$  \\
$0.5$  & $0.368$ & $\phantom{-}0.0$  & $0.316$  & $0.50$  \\
$0.05$ & $0.905$ & $\phantom{-}0.90$ & $0.0475$ & $0.05$  \\
\bottomrule
\end{tabularx}
\end{table}
\end{example}

\begin{pitfall}[思维陷阱：欧拉离散化恒安全]
$T$ 不够小时一阶近似 \cref{eq:ssb-euler} 会\emph{失真甚至失稳}：上表 $T=1$ 时近似的 $A_{d,22}=1-2T=-1$（在单位圆上、临界），而精确值是 $0.135$（稳定）；$T>1$ 时近似值的模超过 $1$，离散系统变\emph{不稳定}，而原连续系统稳定。\emph{欧拉法把稳定系统离散成不稳定，是数字控制的经典翻车点}。所以采样周期要按系统最快模态（最小时间常数）选；要安全无条件用精确 ZOH \cref{eq:ssb-zoh}。这与 \cref{ch:control} 数字 PID 的离散化告诫一脉相承。
\end{pitfall}

\begin{insight}
离散化是连接“连续控制理论”与“真实数字实现”的接口，下游处处用它：MPC（\cref{ch:control} 的滚动时域）需要离散预测方程 $\mathbf{x}_{k+1}=\mathbf{A}_d\mathbf{x}_k+\mathbf{B}_d\mathbf{u}_k$ 来 condensing 成 QP；卡尔曼滤波（\cref{ch:eskf}）的过程模型是离散的；足式的线性倒立摆（LIPM）做 ZOH 离散后才能在线 MPC。\emph{时变离散化}（$\mathbf{A}_d(k)=\boldsymbol{\Phi}((k{+}1)T,kT)$，用 \cref{eq:ssb-peano-baker} 的 Taylor 近似 $\boldsymbol{\Phi}\approx\mathbf{I}+\mathbf{A}(kT)T+\tfrac{T^2}{2}[\mathbf{A}^2+\dot{\mathbf{A}}]$）则服务于沿轨迹的 LTV-MPC。一句话：本章给的 $(\mathbf{A}_d,\mathbf{B}_d)$ 是整个数字/优化控制栈的输入。
\end{insight}

\begin{practice}
\begin{enumerate}
  \item 对 \cref{ex:ssb-discretize} 取 $T=0.1$，用精确 \cref{eq:ssb-zoh} 与近似 \cref{eq:ssb-euler} 各算一遍 $\mathbf{A}_d,\mathbf{B}_d$，比较相对误差。
  \item （增广法）对该例 $\mathbf{A}$ 奇异，用 Van~Loan 增广矩阵 $\exp\!\big(\begin{bsmallmatrix}\mathbf{A}&\mathbf{B}\\\mathbf{0}&\mathbf{0}\end{bsmallmatrix}T\big)$ 验证能一次性得到 \cref{ex:ssb-discretize} 的 $\mathbf{A}_d,\mathbf{B}_d$。
  \item 给定单关节模型 \cref{eq:ssb-rot-ss}（取 $K=0$，即纯惯量$+$阻尼），求其 ZOH 离散化，并讨论采样周期上限。
\end{enumerate}
\end{practice}

% ════════════════════════════════════════════════════════════════
\section{离散状态方程的解：递推法与 Z 变换法}
\label{sec:ssb-discrete-sol}

\paragraph{动机：离散模型自身也要“解出来”。}
上节把连续模型化成离散 $\mathbf{x}(k{+}1)=\mathbf{A}_d\mathbf{x}(k)+\mathbf{B}_d\mathbf{u}(k)$，但只回答了“矩阵从哪来”。真正在数字处理器上跑控制、在 MPC 里写预测方程时，还得回答另一半：给定初值与输入序列，状态序列 $\mathbf{x}(k)$ 究竟是什么？这正是离散版的“求解状态方程”，与连续侧 \cref{sec:ssb-nonhomog} 平行而不可省。为与离散控制文献的习惯记号一致，本节把离散系统记为
\begin{equation}
  \mathbf{x}(k{+}1)=\mathbf{G}\mathbf{x}(k)+\mathbf{H}\mathbf{u}(k),\qquad \mathbf{x}(k)\big|_{k=0}=\mathbf{x}(0),
  \label{eq:ssb-discrete-sys}
\end{equation}
其中 $\mathbf{G}=\mathbf{A}_d=\mathrm{e}^{\mathbf{A}T}$、$\mathbf{H}=\mathbf{B}_d$ 正是上节 \cref{eq:ssb-zoh} 的离散矩阵（也可以是直接给定的离散模型）。离散方程有两种标准解法\cite{liubao2006modern}：\emph{递推法}（亦称迭代法，对定常与时变系统都适用）与 \emph{Z 变换法}（只适用于定常系统，但能一举得到闭式与状态转移矩阵）。

\subsection{递推法：把差分方程一步步迭代}
\label{sec:ssb-discrete-recur}

\begin{theorem}[离散状态方程的递推解]\label{thm:ssb-discrete-recur}
线性定常离散系统 \cref{eq:ssb-discrete-sys} 的解为
\begin{equation}
  \mathbf{x}(k)=\mathbf{G}^k\mathbf{x}(0)+\sum_{j=0}^{k-1}\mathbf{G}^{k-j-1}\mathbf{H}\mathbf{u}(j)
  =\mathbf{G}^k\mathbf{x}(0)+\sum_{j=0}^{k-1}\mathbf{G}^{j}\mathbf{H}\mathbf{u}(k-j-1).
  \label{eq:ssb-discrete-recur}
\end{equation}
\end{theorem}
\begin{proof}
直接对 \cref{eq:ssb-discrete-sys} 自 $k=0$ 起逐步迭代：
\begin{align*}
  k=0:&\quad \mathbf{x}(1)=\mathbf{G}\mathbf{x}(0)+\mathbf{H}\mathbf{u}(0);\\
  k=1:&\quad \mathbf{x}(2)=\mathbf{G}\mathbf{x}(1)+\mathbf{H}\mathbf{u}(1)=\mathbf{G}^2\mathbf{x}(0)+\mathbf{G}\mathbf{H}\mathbf{u}(0)+\mathbf{H}\mathbf{u}(1);\\
  k=2:&\quad \mathbf{x}(3)=\mathbf{G}\mathbf{x}(2)+\mathbf{H}\mathbf{u}(2)=\mathbf{G}^3\mathbf{x}(0)+\mathbf{G}^2\mathbf{H}\mathbf{u}(0)+\mathbf{G}\mathbf{H}\mathbf{u}(1)+\mathbf{H}\mathbf{u}(2);
\end{align*}
归纳到第 $k$ 步得通式
\[
  \mathbf{x}(k)=\mathbf{G}^k\mathbf{x}(0)+\mathbf{G}^{k-1}\mathbf{H}\mathbf{u}(0)+\mathbf{G}^{k-2}\mathbf{H}\mathbf{u}(1)+\dots+\mathbf{G}\mathbf{H}\mathbf{u}(k-2)+\mathbf{H}\mathbf{u}(k-1),
\]
按求和指标整理即 \cref{eq:ssb-discrete-recur} 第一式；把 $j$ 换为 $k-j-1$ 重标即第二式。若自 $k=h$ 起算、初态 $\mathbf{x}(h)$，同理得 $\mathbf{x}(k)=\mathbf{G}^{k-h}\mathbf{x}(h)+\sum_{j=h}^{k-1}\mathbf{G}^{k-j-1}\mathbf{H}\mathbf{u}(j)$。
\end{proof}

\begin{definition}[离散状态转移矩阵]\label{def:ssb-discrete-stm}
对照连续的 $\boldsymbol{\Phi}(t)=\mathrm{e}^{\mathbf{A}t}$，定义离散状态转移矩阵 $\boldsymbol{\Phi}(k):=\mathbf{G}^k$。它满足
\begin{equation}
  \boldsymbol{\Phi}(k{+}1)=\mathbf{G}\boldsymbol{\Phi}(k),\quad \boldsymbol{\Phi}(0)=\mathbf{I},\quad
  \boldsymbol{\Phi}(k-h)=\boldsymbol{\Phi}(k-h_1)\boldsymbol{\Phi}(h_1-h),\quad \boldsymbol{\Phi}^{-1}(k)=\boldsymbol{\Phi}(-k),
  \label{eq:ssb-discrete-stm}
\end{equation}
于是 \cref{eq:ssb-discrete-recur} 写成 $\mathbf{x}(k)=\boldsymbol{\Phi}(k)\mathbf{x}(0)+\sum_{j=0}^{k-1}\boldsymbol{\Phi}(k-j-1)\mathbf{H}\mathbf{u}(j)$，与连续解形式一一对应。
\end{definition}

\begin{insight}
离散解 \cref{eq:ssb-discrete-recur} 与连续解 \cref{eq:ssb-nonhomog} 是“同一首曲子的两个调”：$\mathbf{G}^k$ 对应 $\mathrm{e}^{\mathbf{A}(t-t_0)}$（自由运动），离散卷积和 $\sum\mathbf{G}^{k-j-1}\mathbf{H}\mathbf{u}(j)$ 对应卷积积分 $\int\mathrm{e}^{\mathbf{A}(t-\tau)}\mathbf{B}\mathbf{u}\,\mathrm{d}\tau$（强迫运动）。最关键的\emph{因果性}读法：第 $k$ 步状态只依赖 $\mathbf{u}(0),\dots,\mathbf{u}(k-1)$——当前采样 $\mathbf{u}(k)$ 还来不及影响 $\mathbf{x}(k)$，求和上限恰好停在 $k-1$。这条“当前输入不进当前状态”是离散因果系统的本质，也是后面写 MPC 预测方程时求和指标的由来。
\end{insight}

\paragraph{批量（堆叠）形式：MPC 预测方程的胚胎。}
把 $\mathbf{x}(1),\dots,\mathbf{x}(N)$ 按 \cref{eq:ssb-discrete-recur} 全部展开、叠成一个大向量，得
\begin{equation}
  \begin{bmatrix}\mathbf{x}(1)\\\mathbf{x}(2)\\\vdots\\\mathbf{x}(N)\end{bmatrix}
  =\begin{bmatrix}\mathbf{G}\\\mathbf{G}^2\\\vdots\\\mathbf{G}^N\end{bmatrix}\mathbf{x}(0)
  +\begin{bmatrix}\mathbf{H}&\mathbf{0}&\cdots&\mathbf{0}\\\mathbf{G}\mathbf{H}&\mathbf{H}&\cdots&\mathbf{0}\\\vdots&\vdots&\ddots&\vdots\\\mathbf{G}^{N-1}\mathbf{H}&\mathbf{G}^{N-2}\mathbf{H}&\cdots&\mathbf{H}\end{bmatrix}
  \begin{bmatrix}\mathbf{u}(0)\\\mathbf{u}(1)\\\vdots\\\mathbf{u}(N-1)\end{bmatrix}.
  \label{eq:ssb-discrete-batch}
\end{equation}
这块下三角 Toeplitz 矩阵把“未来状态”写成“初态 $+$ 未来输入序列”的线性函数——正是 MPC 把预测方程 condensing 成二次规划（QP）时的核心结构（\cref{ch:control} 的滚动时域）。换言之，离散递推解不是练习题，而是模型预测控制预测层的代数底座。

\subsection{Z 变换法：一举得到闭式与转移矩阵}
\label{sec:ssb-discrete-ztrans}

递推法直观、可上机，却得不到 $\mathbf{x}(k)$ 关于 $k$ 的\emph{闭式}。对定常系统，Z 变换法能像连续侧的拉氏法那样，把差分方程变成代数方程一举解出。

\begin{theorem}[Z 变换解]\label{thm:ssb-discrete-ztrans}
对 \cref{eq:ssb-discrete-sys} 两端取 Z 变换（用移位性质 $Z[\mathbf{x}(k{+}1)]=z\mathbf{X}(z)-z\mathbf{x}(0)$），得 $(z\mathbf{I}-\mathbf{G})\mathbf{X}(z)=z\mathbf{x}(0)+\mathbf{H}\mathbf{U}(z)$，于是
\begin{equation}
  \mathbf{X}(z)=(z\mathbf{I}-\mathbf{G})^{-1}z\,\mathbf{x}(0)+(z\mathbf{I}-\mathbf{G})^{-1}\mathbf{H}\mathbf{U}(z),
  \label{eq:ssb-discrete-ztrans}
\end{equation}
取反变换 $\mathbf{x}(k)=Z^{-1}[\mathbf{X}(z)]$ 即得解。与递推解 \cref{eq:ssb-discrete-recur} 逐项对照，离散状态转移矩阵有闭式
\begin{equation}
  \mathbf{G}^k=\boldsymbol{\Phi}(k)=Z^{-1}\!\big[(z\mathbf{I}-\mathbf{G})^{-1}z\big].
  \label{eq:ssb-discrete-gk}
\end{equation}
\end{theorem}
\begin{proof}[证（$\mathbf{G}^k$ 的 Z 变换）]
按定义 $Z[\mathbf{G}^k]=\sum_{k=0}^{\infty}\mathbf{G}^k z^{-k}=\mathbf{I}+\mathbf{G}z^{-1}+\mathbf{G}^2z^{-2}+\cdots$。两端左乘 $\mathbf{G}z^{-1}$ 得 $\mathbf{G}z^{-1}Z[\mathbf{G}^k]=\mathbf{G}z^{-1}+\mathbf{G}^2z^{-2}+\cdots$，两式相减消去除首项外所有项，得 $(\mathbf{I}-\mathbf{G}z^{-1})Z[\mathbf{G}^k]=\mathbf{I}$，故 $Z[\mathbf{G}^k]=(\mathbf{I}-\mathbf{G}z^{-1})^{-1}=(z\mathbf{I}-\mathbf{G})^{-1}z$，反变换即 \cref{eq:ssb-discrete-gk}。强迫项用离散卷积定理 $Z\big[\sum_{j=0}^{k-1}\mathbf{G}^{k-j-1}\mathbf{H}\mathbf{u}(j)\big]=(z\mathbf{I}-\mathbf{G})^{-1}\mathbf{H}\mathbf{U}(z)$，与 \cref{eq:ssb-discrete-ztrans} 第二项一致。
\end{proof}

\begin{insight}
\cref{eq:ssb-discrete-gk} 是连续侧拉氏法 \cref{thm:ssb-laplace}（$\mathrm{e}^{\mathbf{A}t}=\mathcal{L}^{-1}[(s\mathbf{I}-\mathbf{A})^{-1}]$）的离散孪生：把 $s\!\leftrightarrow\!z$、$\mathcal{L}^{-1}\!\leftrightarrow\!Z^{-1}$、分母 $(s\mathbf{I}-\mathbf{A})\!\leftrightarrow\!(z\mathbf{I}-\mathbf{G})$，唯一的“离散胎记”是多出一个因子 $z$（因 $Z[\mathbf{G}^k]$ 的级数从 $\mathbf{I}$ 起步）。两侧的极点——连续看 $\mathbf{A}$ 的特征值实部、离散看 $\mathbf{G}$ 的特征值模长——分别决定模态 $\mathrm{e}^{\lambda t}$ 与 $\lambda^k$ 的增减，这正是离散稳定性“特征值落单位圆内”（\cref{thm:lyap-disc}）与连续“落左半平面”对偶的根。
\end{insight}

\begin{example}[递推法与 Z 变换法求同一离散系统]\label{ex:ssb-discrete-sol}
设 $\mathbf{G}=\begin{bmatrix}0&1\\-0.16&-1\end{bmatrix}$、$\mathbf{H}=\begin{bmatrix}1\\1\end{bmatrix}$，初态 $\mathbf{x}(0)=\begin{bmatrix}1\\-1\end{bmatrix}$、单位阶跃输入 $\mathbf{u}(k)=1$\cite{liubao2006modern}。

\noindent \textbf{第一步：转移矩阵 $\boldsymbol{\Phi}(k)=\mathbf{G}^k$.}\;特征方程 $|\,z\mathbf{I}-\mathbf{G}\,|=(z+0.2)(z+0.8)=0$，得 $\lambda_1=-0.2,\ \lambda_2=-0.8$。无论用对角化（$\mathbf{G}=\mathbf{T}\boldsymbol{\Lambda}\mathbf{T}^{-1}$，取 $\mathbf{T}=\begin{bsmallmatrix}1&1\\-0.2&-0.8\end{bsmallmatrix}$）算 $\mathbf{G}^k=\mathbf{T}\,\mathrm{diag}(\lambda_i^k)\,\mathbf{T}^{-1}$，还是用 \cref{eq:ssb-discrete-gk} 对 $(z\mathbf{I}-\mathbf{G})^{-1}z$ 部分分式后取 $Z^{-1}$，都得
\[
  \boldsymbol{\Phi}(k)=\frac{1}{3}\begin{bmatrix}4(-0.2)^k-(-0.8)^k & 5[(-0.2)^k-(-0.8)^k]\\ -0.8[(-0.2)^k-(-0.8)^k] & -(-0.2)^k+4(-0.8)^k\end{bmatrix}.
\]
\noindent \textbf{第二步：完整解 $\mathbf{x}(k)$.}\;Z 变换法最省事：$\mathbf{U}(z)=\tfrac{z}{z-1}$，$z\mathbf{x}(0)+\mathbf{H}\mathbf{U}(z)=\begin{bsmallmatrix}z^2/(z-1)\\(-z^2+2z)/(z-1)\end{bsmallmatrix}$，代入 \cref{eq:ssb-discrete-ztrans} 部分分式后逐元反变换得
\[
  \mathbf{x}(k)=\begin{bmatrix}-\tfrac{17}{6}(-0.2)^k+\tfrac{22}{9}(-0.8)^k+\tfrac{25}{18}\\[4pt] \tfrac{3.4}{6}(-0.2)^k-\tfrac{17.6}{9}(-0.8)^k+\tfrac{7}{18}\end{bmatrix}.
\]
用递推法（对 \cref{eq:ssb-discrete-recur} 的卷积和按等比级数求和）得到\emph{逐字相同}的结果——两法殊途同归。因 $|\lambda_{1,2}|<1$，故 $\lambda_i^k\to0$，状态收敛到由阶跃输入决定的稳态 $\big[\tfrac{25}{18},\,\tfrac{7}{18}\big]^\top$（亦可由 $\mathbf{x}_\infty=(\mathbf{I}-\mathbf{G})^{-1}\mathbf{H}$ 直接核对）。
\end{example}

\begin{insight}
两法分工与连续侧完全平行：\emph{递推法}像连续的数值级数——一步一迭代、对\emph{时变}系统 $\mathbf{x}(k{+}1)=\mathbf{G}(k)\mathbf{x}(k)+\mathbf{H}(k)\mathbf{u}(k)$ 照样能推（只是 $\boldsymbol{\Phi}$ 不再是幂 $\mathbf{G}^k$ 而是连乘 $\mathbf{G}(k{-}1)\cdots\mathbf{G}(0)$，呼应连续 LTV 的基本解阵 \cref{eq:ssb-peano-baker}），且天然就是计算机迭代格式；\emph{Z 变换法}像连续的拉氏法——只对定常系统、但给闭式且顺带暴露极点。工程取舍：在线递推、离线分析用 Z 变换。下游处处接驳：离散李雅普诺夫判据（\cref{thm:lyap-disc}）判 $\mathbf{G}^k$ 是否衰减、卡尔曼滤波（\cref{ch:eskf}）的离散预测、MPC 预测方程 \cref{eq:ssb-discrete-batch} 都直接站在本节之上。
\end{insight}

\begin{practice}
\begin{enumerate}
  \item 对 \cref{ex:ssb-discrete-sol}，用递推法显式写出 $\mathbf{x}(1),\mathbf{x}(2),\mathbf{x}(3)$，并与闭式 $\mathbf{x}(k)$ 在 $k=1,2,3$ 处核对。
  \item 证明离散转移矩阵性质 \cref{eq:ssb-discrete-stm} 的 $\boldsymbol{\Phi}^{-1}(k)=\boldsymbol{\Phi}(-k)$（提示：$\mathbf{G}^k\mathbf{G}^{-k}=\mathbf{I}$，需 $\mathbf{G}$ 非奇异——离散化得到的 $\mathbf{G}=\mathrm{e}^{\mathbf{A}T}$ 恒非奇异）。
  \item （连续—离散对照）取 \cref{ex:ssb-discretize} 的 $\mathbf{A}_d,\mathbf{B}_d$ 作为本节的 $\mathbf{G},\mathbf{H}$，用 Z 变换法写出该系统对单位阶跃的零状态响应 $\mathbf{X}(z)$。
\end{enumerate}
\end{practice}

% ════════════════════════════════════════════════════════════════
\section{本章小结与速查}
\label{sec:ssb-summary}

\paragraph{一条主线回顾。}
本章把“物理系统”一路带到“可在数字处理器上求解的离散模型”：先用三途径（机理/结构图/传函）把电、机、机电系统统一写成 $\dot{\mathbf{x}}=\mathbf{A}\mathbf{x}+\mathbf{B}\mathbf{u}$（\cref{sec:ssb-state,sec:ssb-build}）；用线性变换厘清“矩阵会变、特征值/传函不变”，并以约旦型把系统解耦成模态（\cref{sec:ssb-transform}）；用矩阵指数 $\mathrm{e}^{\mathbf{A}t}$ 把自由运动解出、四算法各取所长（\cref{sec:ssb-expm,sec:ssb-four-algos}）；非齐次解 $=$ 自由 $+$ 卷积强迫（\cref{sec:ssb-nonhomog}）；时变情形用 Peano--Baker/基本解阵定义 $\boldsymbol{\Phi}(t,t_0)$，为机器人沿轨迹线性化打底（\cref{sec:ssb-ltv}）；最后 ZOH 离散化供给 $(\mathbf{A}_d,\mathbf{B}_d)$（\cref{sec:ssb-discretize}）。

\begin{table}[H]\centering\small
\caption{核心符号速查}
\label{tab:ssb-symbols}
\begin{tabular}{lll}
\toprule
符号 & 含义 & 出处 \\
\midrule
$\mathbf{x}\in\mathbb{R}^n,\ \mathbf{u}\in\mathbb{R}^r,\ \mathbf{y}\in\mathbb{R}^m$ & 状态/输入/输出矢量 & \cref{def:ssb-state} \\
$\mathbf{A},\mathbf{B},\mathbf{C},\mathbf{D}$ & 系统/输入/输出/直传矩阵 & \cref{def:ssb-ss} \\
$\boldsymbol{\Phi}(t)=\mathrm{e}^{\mathbf{A}t}$ & 定常状态转移矩阵 & \cref{def:ssb-stm} \\
$\boldsymbol{\Phi}(t,t_0)$ & 时变状态转移矩阵（两时刻）& \cref{def:ssb-ltv-stm} \\
$\mathbf{T}$（非奇异）& 坐标变换阵 $\mathbf{x}=\mathbf{T}\mathbf{z}$ & \cref{eq:ssb-similarity} \\
$\boldsymbol{\Lambda},\mathbf{J}$ & 对角型/约旦型 & \cref{sec:ssb-jordan} \\
$\mathbf{A}_d=\mathrm{e}^{\mathbf{A}T},\ \mathbf{B}_d$ & ZOH 离散化矩阵 & \cref{eq:ssb-zoh} \\
\bottomrule
\end{tabular}
\end{table}

\begin{table}[H]\centering\small
\caption{定理与公式速查}
\label{tab:ssb-theorems}
\begin{tabular}{lll}
\toprule
结论 & 内容 & 标号 \\
\midrule
不变量 & 特征值/特征多项式系数在相似变换下不变 & \cref{thm:ssb-invariant} \\
矩阵指数 & $\mathrm{e}^{\mathbf{A}t}=\sum_k\tfrac{1}{k!}\mathbf{A}^k t^k$；$\mathrm{e}^{\mathbf{A}t}\mathrm{e}^{\mathbf{B}t}{=}\mathrm{e}^{(\mathbf{A}+\mathbf{B})t}{\iff}\mathbf{A}\mathbf{B}{=}\mathbf{B}\mathbf{A}$ & \cref{eq:ssb-expm-def,eq:ssb-stm-commute} \\
四算法 & 级数 / 约旦型 / 拉氏反变换 / Cayley--Hamilton & \cref{sec:ssb-four-algos} \\
非齐次解 & $\mathbf{x}=\mathrm{e}^{\mathbf{A}(t-t_0)}\mathbf{x}_0+\int\mathrm{e}^{\mathbf{A}(t-\tau)}\mathbf{B}\mathbf{u}\,\mathrm{d}\tau$ & \cref{thm:ssb-nonhomog} \\
时变解 & $\boldsymbol{\Phi}(t,t_0)$ 由 Peano--Baker 级数/基本解阵 & \cref{thm:ssb-peano-baker} \\
离散化 & $\mathbf{A}_d=\mathrm{e}^{\mathbf{A}T},\ \mathbf{B}_d=(\int_0^T\mathrm{e}^{\mathbf{A}t}\mathrm{d}t)\mathbf{B}$ & \cref{thm:ssb-zoh} \\
\bottomrule
\end{tabular}
\end{table}

\subsection*{常见误解汇总}
\begin{itemize}
  \item “状态空间表达式是唯一的”——错，状态选取不唯一（\cref{sec:ssb-invariant}），但特征值/传函是不变量。
  \item “$\mathrm{e}^{\mathbf{A}t}\mathrm{e}^{\mathbf{B}t}=\mathrm{e}^{(\mathbf{A}+\mathbf{B})t}$ 总成立”——错，要 $\mathbf{A}\mathbf{B}=\mathbf{B}\mathbf{A}$（\cref{eq:ssb-stm-commute}）。
  \item “时变系统也能写 $\mathrm{e}^{\int\mathbf{A}\mathrm{d}\tau}$”——一般不能，须可交换条件（\cref{thm:ssb-ltv-commute}），否则用 Peano--Baker。
  \item “$m\neq0$ 时直接拿 $y,\dot y,\dots$ 当状态”——会引入 $\dot u$ 致解病态（\cref{eq:ssb-bad-state}），须用能控/能观标准形。
  \item “欧拉离散化总是安全的”——$T$ 大会失稳（\cref{ex:ssb-discretize} 表），要看最快模态或用精确 ZOH。
\end{itemize}

\subsection*{通向高级}
本章建立的“统一状态空间记号底座 $+$ 求解工具”是 C0 子部与 A--Q 高级章的共同地基。向下游延伸：
\begin{itemize}
  \item \textbf{能控能观}（\cref{ch:controllability-obsv}）：本章的非齐次解 \cref{eq:ssb-nonhomog} 与约旦模态视角，正是能控性（输入能否激发各模态）、能观性（输出能否看到各模态）的出发点；本章给的能控/能观标准形（\cref{der:ssb-canonical}）将在那里讲透其命名由来与最小实现。
  \item \textbf{稳定性}（\cref{ch:lyapunov-basics}）：$\mathrm{e}^{\mathbf{A}t}$ 的模态 $\mathrm{e}^{\lambda_i t}$ 直接决定渐近稳定（$\mathrm{Re}\,\lambda_i<0$），是李雅普诺夫第一法（间接法）与李氏方程 $\mathbf{A}^\top\mathbf{P}+\mathbf{P}\mathbf{A}=-\mathbf{Q}$（\cref{eq:ctrl-lyap-eq}）的根。
  \item \textbf{综合与最优控制}（\cref{ch:linear-synthesis,ch:optimal-control-basics}）：状态反馈 $\mathbf{u}=-\mathbf{K}\mathbf{x}$ 改变闭环 $\mathbf{A}-\mathbf{B}\mathbf{K}$ 的谱、观测器重构状态、LQR 在 \cref{eq:ssb-nonhomog} 的解上优化输入——全部以本章的 $(\mathbf{A},\mathbf{B})$ 与 $\mathrm{e}^{\mathbf{A}t}$ 为前提。
  \item \textbf{动力学与数字/优化控制}：\cref{eq:ssb-mechanical-general} 把 \cref{ch:rigid_body} 的二阶动力学降阶为状态空间，供 C2 动力学（$\mathbf{J}/\mathbf{K}/\mathbf{L}$）取用；ZOH 离散化 \cref{eq:ssb-zoh} 为 MPC condensing、LIPM 在线规划、卡尔曼滤波（\cref{ch:eskf}）供给离散 $(\mathbf{A}_d,\mathbf{B}_d)$；LTV（\cref{sec:ssb-ltv}）为 iLQR/DDP 沿轨迹线性化打底。
\end{itemize}
读到这里，你已握有“把任意（线性化的）机器人系统写成状态空间、解出其演化、并离散化交给计算机”的全套地基工具。
